% ================================================================
%  第 5 章  极小极大构造
% ================================================================
\chapter{极小极大构造}
\label{chap:5}

{\setlength{\emergencystretch}{2em}%
本章使用一个非常一般的论证来在任意球面上寻找极小球面, 其基本想法可追溯到
H. A. Schwarz 1870 年代关于调和函数的工作, 以及 G. D. Birkhoff 1910 年代
关于闭测地线的工作. 这里的处理遵循文献 \cite{CM27} 与 \cite{CM28}, 其中
若干存在性结果是新的.\par}

Schwarz 与 Birkhoff 的想法都是用一个极小极大的论证来证明变分问题的临界点
的存在性. 在本章中, 这一想法使我们尤其能够构造出非稳定的极小曲面. 在极小
曲面的极小极大构造中, 人们用一族曲面扫掠流形, 同时记录扫掠中各切片的面积;
然后试图抽出一个收敛的极大切片序列, 使这些极大切片的面积收敛到所有扫掠的
极大切片面积的下确界. 其想法是: 只要做得恰当, 这样的极大切片序列就应当在
某种弱拓扑下收敛到一张非平凡的极小曲面. 尽管这种构造常被称为极小极大构造,
它实际上应当称为下确界--极大构造, 因为所取的是所有扫掠的极大面积的下确界.

% ================================================================
\section{用曲线扫掠}
\label{sec:5.1}

给定 $2$-球面上的一个 Riemann 度量, 用一族连续的、起止于点曲线的闭曲线
来扫掠该 $2$-球面. 再以连续的方式把每条曲线尽可能拉紧、同时保持扫掠不变,
从而把这个扫掠收紧.

用扫掠来寻找 $2$-球面上的闭测地线可追溯到 Birkhoff 1917 年的工作; 关于
Birkhoff 想法的更多内容见 \cite{B1}、\cite{B2} 以及 \cite{Cr} 的第 2 节.
该论证在任意闭流形上同样有效, 但只有当宽度(其定义见下文 \eqref{eq:5.23})
为正时才能给出非平凡的闭测地线. 例如, 当 $M$ 在拓扑上是一个 $2$-球面时,
粗略地说, 宽度在相差一个常数的意义下等于把 $M$ ``拉过来''所需的最短闭
曲线长度的平方. Birkhoff 的论证给出: $2\pi$ 乘以宽度可实现为某条闭测地线
长度的平方.

本章中 $M$ 总是一个闭 Riemann 流形, 等距嵌入某个欧氏空间 $\R^N$. 为方便
起见, 我们用常数缩放 $\R^N$(从而缩放 $M$), 使其满足下述三条性质:

\begin{enumerate}[label=(M\arabic*),leftmargin=3.2em]
\item $\sup_M|A|\leq 1/16$, 其中 $|A|^2$ 是 $M$ 的第二基本形式的范数平方;
\item $M$ 的单射半径至少为 $8\pi$, 且曲率至多为 $1/64$, 于是 $M$ 中半径
      不超过 $4\pi$ 的每个测地球都严格测地凸;
\item 若 $x,y\in M$ 且 $|x-y|\leq 1$, 则 $\dist_M(x,y)\leq 2|x-y|$.
\end{enumerate}

最后一条性质是联系内在距离与外在距离的弦弧界.

% ================================================================
\section{Birkhoff 的曲线缩短过程}
\label{sec:5.2}

记 $W^{1,2}$ 为从 $\bS^1$ 到 $M$ 的 $W^{1,2}$ 映射构成的空间, 其中把
$\bS^1$ 与 $\R/2\pi\mathbb{Z}$ 等同. 我们在 $W^{1,2}$ 上使用由
$W^{1,2}$(Sobolev) 范数给出的距离与拓扑. 回忆映射 $f:\bS^1\to\R^N$ 的
$W^{1,2}$ 范数的平方为
\begin{equation*}
\int_{\bS^1}\bigl(|f|^2+|f'|^2\bigr).
\end{equation*}

关于 $W^{1,2}$ 范数下的曲线空间, 我们将用到两条基本事实.

第一条基本事实是: $W^{1,2}$ 范数蕴含 H\"older 正则性. 为此设
$x,y\in\bS^1$, 用微积分基本定理与 Cauchy--Schwarz 不等式可得
\begin{equation}
\begin{aligned}
|f(x)-f(y)|^2
&\leq\Bigl(\int_x^y|f'|\Bigr)^2\leq\Bigl(\int_x^y 1\Bigr)
 \Bigl(\int_x^y|f'|^2\Bigr)\\
&\leq|x-y|\int_{\bS^1}|f'|^2.
\end{aligned}
\label{eq:5.1}
\end{equation}
于是 $f\in C^{1/2}$. 因此, 若 $f,g$ 在 $W^{1,2}$ 意义下接近, 则把上式
用于 $(f-g)$ 便知它们在 $C^0$ 意义下接近.

{\setlength{\emergencystretch}{2em}%
我们将反复使用的第二条基本估计是 Wirtinger 不等式. 它不过是通常的
Dirichlet--Poincar\'e 不等式, 即用导数的 $L^2$ 范数来控制 $L^2$ 范数:\par}
\begin{equation*}
\int_0^{\pi L}f^2\,dt\leq L^2\int_0^{\pi L}(f')^2\,dt,
\end{equation*}
只要 $f(0)=f(\pi L)=0$. 注意当 $f(t)=a\sin(t/L)$($a$ 为常数)时 Wirtinger
不等式取等号.

\mtpara{曲线空间与曲线缩短过程}
固定一个大的正整数 $L$, 记 $\Lambda$ 为从 $\bS^1$ 到 $M$ 的、恰有 $L$ 个
折点的分段线性映射构成的空间(允许有冗余的折点), 其中每段测地线的长度
至多为 $2\pi$, 按弧长的(常数)倍数参数化, 并具有 Lipschitz 界 $L$. 所谓
线性映射, 我们指的是(常速)测地线. 记 $G\subset\Lambda$ 为 $M$ 中长度
至多 $2\pi L$ 的浸入闭测地线之集. ($\Lambda$ 中曲线的能量等于其长度的
平方除以 $2\pi$. 换言之, 能量与长度本质上是等价的.) 注意 $\Lambda$ 是
有限维的, 其维数由 $L$ 与 $M$ 的维数给出.

我们在 $\Lambda$ 上使用由从 $\bS^1$ 到 $M$ 的映射空间上的 $W^{1,2}$ 范数
给出的距离与拓扑.

曲线缩短是一个映射 $\Psi:\Lambda\to\Lambda$, 使得\footnote{这个映射本质上
就是所谓的 Birkhoff 曲线缩短过程; 见 \cite{Cr} 第 2 节或 \cite{CM28}.}:

\begin{enumerate}
\item $\Psi(\gamma)$ 与 $\gamma$ 同伦, 且
      $\Length(\Psi(\gamma))\leq\Length(\gamma)$.
\item $\Psi(\gamma)$ 连续依赖于 $\gamma$.
\item 存在连续函数 $\phi:[0,\infty)\to[0,\infty)$, $\phi(0)=0$, 使得
\begin{equation*}
\dist^2(\gamma,\Psi(\gamma))
\leq\phi\Bigl(\frac{\Length^2(\gamma)-\Length^2(\Psi(\gamma))}
{\Length^2(\Psi(\gamma))}\Bigr).
\end{equation*}
\item 给定 $\epsilon>0$, 存在 $\delta>0$, 使得若 $\gamma\in\Lambda$ 且
      $\dist(\gamma,G)\geq\epsilon$, 则
      $\Length(\Psi(\gamma))\leq\Length(\gamma)-\delta$.
\end{enumerate}

这些性质的若干版本已为 Birkhoff 所知, 但我们这里给出的强形式是在
\cite{CM28} 中证明的.

\mtpara{定义曲线缩短 $\Psi$}
为定义 $\Psi$, 我们固定 $\bS^1$ 的一个分划: 取 $2L$ 个依次等距排列的点%
\footnote{注意这些点未必是分段线性映射的折点所在之处.}
\begin{equation}
x_0,x_1,x_2,\dots,x_{2L}=x_0\in\bS^1,
\label{eq:5.2}
\end{equation}
使得 $|x_j-x_{j+1}|=\pi/L$. $\Psi(\gamma)$ 分三步给出. 首先对 $\gamma$
施行第一步得到曲线 $\gamma_e$, 再对 $\gamma_e$ 施行第二步得到曲线
$\gamma_o$. 第三步重新参数化 $\gamma_o$, 得到 $\Psi(\gamma)$.

\textbf{第一步}: 把 $\gamma$ 在每个\emph{偶}区间, 即 $[x_{2j},x_{2j+2}]$,
上替换为具有相同端点的线性映射, 得到分段线性曲线 $\gamma_e:\bS^1\to M$.
也就是说, 对每个 $j$, 令 $\gamma_e|_{[x_{2j},x_{2j+2}]}$ 为从
$\gamma(x_{2j})$ 到 $\gamma(x_{2j+2})$ 的唯一最短(常速)测地线.

\textbf{第二步}: 把 $\gamma_e$ 在每个\emph{奇}区间上替换为具有相同端点的
线性映射, 得到分段线性曲线 $\gamma_o:\bS^1\to M$.

\textbf{第三步}: 重新参数化 $\gamma_o$(固定 $\gamma_o(x_0)$), 得到所需的
常速曲线 $\Psi(\gamma):\bS^1\to M$.

容易看出 $\Psi$ 把 $\Lambda$ 映入 $\Lambda$ 并具有性质 (1); 参见 \cite{Cr}
第 2 节.

为证明 (2) 与 (3), 宜于注意: 存在一种等价但更对称的方式, 用四步来构造
$\Psi(\gamma)$:

\begin{itemize}[leftmargin=2.6em,itemsep=1pt,topsep=2pt]
\item[$(A_1)$] 按第一步得到 $\gamma_e$.
\item[$(B_1)$] 重新参数化 $\gamma_e$(固定 $x_0$ 的像), 得到常速曲线
      $\tilde{\gamma}_e$. 此重新参数化把点 $x_j$ 移到新点 $\tilde{x}_j$
      (即 $\gamma_e(x_j)=\tilde{\gamma}_e(\tilde{x}_j)$).
\item[$(A_2)$] 在奇的 $\tilde{x}_j$ 区间上作线性替换, 得到
      $\tilde{\gamma}_o$.
\item[$(B_2)$] 重新参数化 $\tilde{\gamma}_o$(固定 $x_0$ 的像), 得到常速
      曲线 $\Psi(\gamma)$.
\end{itemize}

之所以得到同一条曲线, 是因为 $\tilde{\gamma}_o$ 只是 $\gamma_o$ 的重新
参数化. 我们还将用到: 这四步中的每一步都不增加能量. 对线性替换这是显然的,
因为线性映射使能量极小. 对重新参数化, 它由 Cauchy--Schwarz 不等式推出,
因为对曲线 $\sigma:\bS^1\to M$ 有
\begin{equation}
\Length^2(\sigma)\leq 2\pi\,\mathrm{E}(\sigma),
\label{eq:5.3}
\end{equation}
且等号成立当且仅当 $|\sigma'|=\Length(\sigma)/(2\pi)$ 几乎处处成立.

\mtpara{$\Psi$ 的性质 (3)}
利用上述四步定义 $\Psi(\gamma)$ 的替代方式可以看出: 只要用长度的减少量
来控制 $\dist(\gamma,\gamma_e)$ 与 $\dist(\gamma_e,\tilde{\gamma}_e)$
(以及步骤 $(A_2)$、$(B_2)$ 的相应量), 性质 (3) 便由三角不等式推出.

对 $\dist(\gamma,\gamma_e)$ 的界直接来自下述引理:

\begin{lemma}
\label{lem:5.1}
存在 $C$, 使得若 $I$ 是长度至多 $2\pi/L$ 的区间, $\sigma_1:I\to M$ 是
满足 $|\sigma_1'|\leq L$ 的 Lipschitz 曲线, 而 $\sigma_2:I\to M$ 是具有
相同端点的极小测地线, 则
\begin{equation}
\dist^2(\sigma_1,\sigma_2)\leq C\bigl(\mathrm{E}(\sigma_1)
-\mathrm{E}(\sigma_2)\bigr).
\label{eq:5.4}
\end{equation}
\end{lemma}

我们将需要 (M1) 与 (M3) 的一个简单推论:

\begin{lemma}
\label{lem:5.2}
若 $x,y\in M$, 则 $|(x-y)^{\perp}|\leq|x-y|^2$, 其中 $(x-y)^{\perp}$ 是
$x-y$ 在 $y$ 处关于 $M$ 的法向分量.
\end{lemma}

\begin{proof}
若 $|x-y|\geq 1$, 则结论显然. 因此设
\begin{equation*}
|x-y|<1,
\end{equation*}
并设 $\alpha:[0,\ell]\to M$ 是从 $y$ 到 $x$ 的极小化单位速度测地线, 且
$\ell\leq 2|x-y|$. 令 $V$ 为单位法向量
\begin{equation*}
V=\frac{(x-y)^{\perp}}{|(x-y)^{\perp}|},
\end{equation*}
于是 $\langle\alpha'(0),V\rangle=0$. 注意到
\begin{equation}
\begin{aligned}
|(x-y)^{\perp}|
&=\int_0^{\ell}\langle\alpha'(s),V\rangle\,ds
 =\int_0^{\ell}\Bigl\langle\alpha'(0)+\int_0^s\alpha''(t)\,dt,V\Bigr\rangle ds\\
&\leq\int_0^{\ell}\int_0^s|\alpha''(t)|\,dt\,ds
 \leq\int_0^{\ell}\int_0^s|A(\alpha(t))|\,dt\,ds\\
&\leq\frac12\ell^2\sup_M|A|\leq|x-y|^2.
\end{aligned}
\label{eq:5.5}
\end{equation}
\end{proof}

\begin{proof}[引理~\ref{lem:5.1}的证明]
分部积分并利用 $\sigma_1$ 与 $\sigma_2$ 在 $\partial I$ 上相等, 得
\begin{equation}
\int_I|\sigma_1'|^2-\int_I|\sigma_2'|^2
-\int_I|(\sigma_1-\sigma_2)'|^2
=-2\int_I\bigl\langle(\sigma_1-\sigma_2),\sigma_2''\bigr\rangle
\equiv\kappa.
\label{eq:5.6}
\end{equation}
只要用 $\frac12\int_I|(\sigma_1-\sigma_2)'|^2$ 来界定 $|\kappa|$, 再诉诸
Wirtinger 不等式, 引理便得证.

由于 $\sigma_2$ 是 $M$ 上的测地线, $\sigma_2''$ 垂直于 $M$, 且
$|\sigma_2''|\leq|\sigma_2'|^2\sup_M|A|\leq|\sigma_2'|^2/16$. 于是
引理~\ref{lem:5.2} 给出
\begin{equation}
\bigl|\bigl\langle(\sigma_1-\sigma_2),\sigma_2''\bigr\rangle\bigr|
\leq|(\sigma_1-\sigma_2)^{\perp}|\,\frac{|\sigma_2'|^2}{16}
\leq|\sigma_1-\sigma_2|^2\,\frac{|\sigma_2'|^2}{16}.
\label{eq:5.7}
\end{equation}
对 \eqref{eq:5.7} 积分, 利用 $|\sigma_2'|$ 为常数且
$|\sigma_2'|\Length(I)\leq 2\pi$, 再应用 Wirtinger 不等式, 得
\begin{equation}
\begin{aligned}
|\kappa|
&\leq\frac{|\sigma_2'|^2}{8}\int_I|\sigma_1-\sigma_2|^2
 \leq\frac{|\sigma_2'|^2}{8}
 \Bigl(\frac{\Length(I)}{\pi}\Bigr)^2\int_I|(\sigma_1-\sigma_2)'|^2\\
&\leq\frac12\int_I|(\sigma_1-\sigma_2)'|^2.
\end{aligned}
\label{eq:5.8}
\end{equation}
\end{proof}

对步骤 $(A_1)$ 中的 $L$ 个区间分别应用引理~\ref{lem:5.1}, 给出
\begin{equation}
\begin{aligned}
\dist^2(\gamma,\gamma_e)
&\leq C\bigl(\mathrm{E}(\gamma)-\mathrm{E}(\gamma_e)\bigr)\\
&\leq\frac{C}{2\pi}\bigl(\Length^2(\gamma)-\Length^2(\Psi(\gamma))\bigr).
\end{aligned}
\label{eq:5.9}
\end{equation}
由于 $\Length(\Psi(\gamma))\leq 2\pi L$, 这就给出了所需的
$\dist(\gamma,\gamma_e)$ 的界.

为界定 $\dist(\gamma_e,\tilde{\gamma}_e)$, 我们将用到 $\gamma_e$ 恰是
复合 $\tilde{\gamma}_e\circ P$, 其中 $P:\bS^1\to\bS^1$ 是单调的分段线性
映射.\footnote{映射 $P$ 是 Lipschitz 的, 但若 $\gamma_e$ 在某个区间上为
常数, 其逆映射 $P^{-1}$ 未必如此.} 利用
$|\tilde{\gamma}_e'|=\Length(\tilde{\gamma}_e)/(2\pi)$(在折点之外)以及
$P'$ 的积分为 $2\pi$, 容易算出
\begin{equation}
\begin{aligned}
\int(P'-1)^2
&=\int(P')^2-2\pi=\int\Bigl(\frac{|\gamma_e'|}
{|\tilde{\gamma}_e'\circ P|}\Bigr)^2-2\pi\\
&=\frac{4\pi^2}{\Length^2(\tilde{\gamma}_e)}\int|\gamma_e'|^2-2\pi
 =2\pi\,\frac{\mathrm{E}(\gamma_e)-\mathrm{E}(\tilde{\gamma}_e)}
{\mathrm{E}(\tilde{\gamma}_e)}\\
&\leq 2\pi\,\frac{\mathrm{E}(\gamma)-\mathrm{E}(\Psi(\gamma))}
{\mathrm{E}(\Psi(\gamma))}.
\end{aligned}
\label{eq:5.10}
\end{equation}
由于 $\gamma_e$ 与 $\tilde{\gamma}_e$ 在 $x_0=x_{2L}$ 处相等, Wirtinger
不等式把 $\dist^2(\gamma_e,\tilde{\gamma}_e)$ 界定为
\begin{equation}
\frac12\int\bigl|(\tilde{\gamma}_e\circ P)'-\tilde{\gamma}_e'\bigr|^2
\leq\int\bigl|(\tilde{\gamma}_e'\circ P)P'-\tilde{\gamma}_e'\circ P\bigr|^2
+\bigl|\tilde{\gamma}_e'\circ P-\tilde{\gamma}_e'\bigr|^2.
\label{eq:5.11}
\end{equation}
我们将用 $\int|P'-1|^2$ 来界定 \eqref{eq:5.11} 右端的两项, 然后诉诸
\eqref{eq:5.10}. 为界定第一项, 利用 $|\tilde{\gamma}_e'|$(为常数)
$\leq L$, 得
\begin{equation}
\int\bigl|(\tilde{\gamma}_e'\circ P)P'-\tilde{\gamma}_e'\circ P\bigr|^2
\leq L^2\int|P'-1|^2.
\label{eq:5.12}
\end{equation}
为界定第二个积分, 我们将用到: 当 $x,y$ 是 $\bS^1$ 中不被折点分隔的两点时,
$\tilde{\gamma}_e$ 是从 $x$ 到 $y$ 的测地线, 因而 $\tilde{\gamma}_e''$
垂直于 $M$, 且
\begin{equation*}
|\tilde{\gamma}_e''|\leq|\tilde{\gamma}_e'|^2\sup_M|A|\leq\frac{L^2}{16}.
\end{equation*}
因此, 把 $\tilde{\gamma}_e''$ 从 $x$ 积分到 $y$ 给出
\begin{equation}
|\tilde{\gamma}_e'(x)-\tilde{\gamma}_e'(y)|
\leq|x-y|\sup|\tilde{\gamma}_e''|\leq\frac{L^2}{16}|x-y|.
\label{eq:5.13}
\end{equation}

把 $\bS^1$ 分为两个集合 $S_1$ 与 $S_2$, 其中 $S_1$ 是到
$\tilde{\gamma}_e$ 的某个折点距离不超过
$\bigl(\pi\int|P'-1|^2\bigr)^{1/2}$ 的点集. 由于 $P(x_0)=x_0$, 仿照
\eqref{eq:5.1} 的论证给出
\begin{equation*}
|P(x)-x|\leq\Bigl(\pi\int|P'-1|^2\Bigr)^{1/2}.
\end{equation*}
于是, 若 $x\in S_2$, 则 $\tilde{\gamma}_e$ 在 $x$ 与 $P(x)$ 之间是光滑的.
从而 \eqref{eq:5.13} 给出
\begin{equation}
\int_{S_2}\bigl|\tilde{\gamma}_e'\circ P-\tilde{\gamma}_e'\bigr|^2
\leq\frac{L^4}{256}\int_{S_2}|P(s)-s|^2
\leq\frac{L^4}{64}\int|P'-1|^2,
\label{eq:5.14}
\end{equation}
其中最后一个不等式用了 Wirtinger 不等式. 另一方面,
\begin{equation}
\int_{S_1}\bigl|\tilde{\gamma}_e'\circ P-\tilde{\gamma}_e'\bigr|^2
\leq 4L^2\Length(S_1)
\leq 8L^3\Bigl(\pi\int|P'-1|^2\Bigr)^{1/2},
\label{eq:5.15}
\end{equation}
这就完成了性质 (3) 的证明.

\mtpara{性质 (4)}
为证明性质 (4), 我们用反证法. 设存在 $\epsilon>0$ 以及序列
$\gamma_j\in\Lambda$, 使得
\begin{equation}
\mathrm{E}(\Psi(\gamma_j))\geq\mathrm{E}(\gamma_j)-1/j,
\label{eq:5.16}
\end{equation}
\begin{equation}
\dist(\gamma_j,G)\geq\epsilon>0.
\label{eq:5.17}
\end{equation}
注意第二个条件蕴含着 $\mathrm{E}(\gamma_j)$ 有正的下界. 其次注意空间
$\Lambda$ 是紧的, 因为 $\sigma\in\Lambda$ 连续依赖于 $L$ 个折点在紧流形
$M$ 中的像. 因此, $\gamma_j$ 的某个子列必须收敛到某个
$\gamma\in\Lambda$. 由于性质 (3) 蕴含
\begin{equation*}
\dist\bigl(\gamma_j,\Psi(\gamma_j)\bigr)\to 0,
\end{equation*}
所以 $\Psi(\gamma_j)$ 也收敛到 $\gamma$. 于是 $\Psi$ 的连续性, 即 $\Psi$
的性质 (2), 蕴含 $\Psi(\gamma)=\gamma$. 然而由于 $\Psi$ 的不动点只能是
浸入闭测地线, 这意味着 $\gamma\in G$. 这后一事实 Birkhoff 已经用过
(见 \cite{Cr} 第 2 节), 它可由引理~\ref{lem:5.1} 与 \eqref{eq:5.10}
立即推出. 但这与 $\gamma_j$ 与任何这样的闭浸入测地线保持固定距离相矛盾,
从而完成了 (4) 的证明.

\mtparasec{性质 (2): $\Psi$ 的连续性}

\begin{lemma}
\label{lem:5.3}
设 $\gamma:\bS^1\to M$ 是满足 $\mathrm{E}(\gamma)\leq L$ 的 $W^{1,2}$
映射. 若 $\gamma_e$ 与 $\tilde{\gamma}_e$ 由对 $\gamma$ 施行步骤
$(A_1)$ 与 $(B_1)$ 得到, 则映射 $\gamma\mapsto\tilde{\gamma}_e$ 是从
$W^{1,2}$ 到赋以 $W^{1,2}$ 范数的 $\Lambda$ 的连续映射.
\end{lemma}

\begin{proof}
由 \eqref{eq:5.1} 与能量界推出
\begin{equation}
\dist_M\bigl(\gamma(x_{2j}),\gamma(x_{2j+2})\bigr)\leq 2\pi
\label{eq:5.18}
\end{equation}
对每个 $j$ 成立, 因此可以施行步骤 $(A_1)$. 引理将由下述两条观察推出:

(C1) 由于 $W^{1,2}$ 接近的曲线也 $C^0$ 接近, 点
$\gamma_e(x_{2j})=\gamma(x_{2j})$ 关于 $W^{1,2}$ 范数是连续的.

(C2) 定义 $\Gamma\subset M\times M$ 为
\begin{equation*}
\Gamma=\{(x,y)\in M\times M\mid \dist_M(x,y)\leq 4\pi\},
\end{equation*}
并定义映射
\begin{equation*}
H:\Gamma\to C^1([0,1],M),
\end{equation*}
其中 $H(x,y):[0,1]\to M$ 取为从 $x$ 到 $y$ 的线性映射. 则 $H$ 在
$\Gamma$ 上连续. 此外, 映射 $t\mapsto H(x,y)(t)$ 具有一致有界的一阶与
二阶导数
\begin{equation*}
|\partial_t H(x,y)|\leq 4\pi,\qquad
|\partial_t^2 H(x,y)|\leq\pi^2.
\end{equation*}
二阶导数的界来自 (M1).\footnote{(M1) 在本章开头定义.}

为证明引理, 设 $\gamma^1$ 与 $\gamma^2$ 是 $\Lambda$ 中的非常值曲线
(在常值映射处的连续性显然). 对 $i=1,2$ 与 $j=1,\dots,L$, 记 $a_j^i$ 为
$M$ 中从 $\gamma^i(x_{2j})$ 到 $\gamma^i(x_{2j+2})$ 的距离. 令
$S^i=\frac1{2\pi}\sum_{j=1}^La_j^i$ 为 $\tilde{\gamma}_e^i$ 的速度,
于是在 $L$ 个折点之外有 $|(\tilde{\gamma}_e^i)'|=S^i$. 由 (C1),
$a_j^i$ 是 $\gamma^i$ 的连续函数, 因而 $S^1$ 与 $S^2$ 也是. 此外,
(C1) 与 (C2) 蕴含 $\gamma_e^1$ 与 $\gamma_e^2$ 在每个区间
$[x_{2j},x_{2j+2}]$ 上 $C^1$ 接近. 这样便证明了 $\gamma\mapsto\gamma_e$
连续.

为证明 $\gamma_e\mapsto\tilde{\gamma}_e$ 也连续, 我们将证明当
$\gamma_e^i$ 接近时 $\tilde{\gamma}_e^i$ 也接近. 由于点 $x_0=x_{2L}$ 在
重新参数化下不动, 只要证明 $\int_{\bS^1}|(\tilde{\gamma}_e^1
-\tilde{\gamma}_e^2)'|^2$ 可以变得任意小, 把 Wirtinger 不等式用于
$(\tilde{\gamma}_e^1-\tilde{\gamma}_e^2)
-(\tilde{\gamma}_e^1-\tilde{\gamma}_e^2)(x_0)$ 即可.

分段线性曲线 $\tilde{\gamma}_e^i$ 在区间
\begin{equation}
I_j^i=\Bigl[\frac1{S^i}\sum_{\ell<j}a_\ell^i,\;
\frac1{S^i}\sum_{\ell\leq j}a_\ell^i\Bigr]
\label{eq:5.19}
\end{equation}
上是线性的. 令 $I_j=I_j^1\cap I_j^2$. 首先注意: 由于 \eqref{eq:5.19} 中的
区间 $I_j^i$ 连续依赖于 $\gamma_e^i$, 补集
$\bS^1\setminus\bigl[\bigcup_{j=1}^LI_j\bigr]$ 的测度可以变得任意小, 因此
\begin{equation}
\int_{\bS^1\setminus[\bigcup I_j]}
\bigl|(\tilde{\gamma}_e^1-\tilde{\gamma}_e^2)'\bigr|^2
\leq 4L^2\Length\bigl(\bS^1\setminus\bigl[\textstyle\bigcup I_j\bigr]\bigr)
\label{eq:5.20}
\end{equation}
也可以变得任意小. 我们将按 $a_j^1$ 的大小把 $I_j$ 分成两组. 固定
$\epsilon>0$, 先设 $a_j^1<\epsilon$; 由连续性可设 $a_j^2<2\epsilon$. 对这样
的 $j$ 得到
\begin{equation}
\begin{aligned}
\int_{I_j}\bigl|(\tilde{\gamma}_e^1-\tilde{\gamma}_e^2)'\bigr|^2
&\leq 2\int_{I_j^1}\bigl|(\tilde{\gamma}_e^1)'\bigr|^2
+2\int_{I_j^2}\bigl|(\tilde{\gamma}_e^2)'\bigr|^2\\
&\leq 2L(a_j^1+a_j^2)\leq 6\epsilon L.
\end{aligned}
\label{eq:5.21}
\end{equation}
由于折点至多有 $L$ 个, 对这些区间求和对
$(\tilde{\gamma}_e^1-\tilde{\gamma}_e^2)$ 的能量至多贡献 $6\epsilon L^2$.

剩下要考虑的是 $a_j^1\geq\epsilon$ 的 $I_j$; 由连续性可设
$a_j^2\geq\epsilon/2$. 此时 $\tilde{\gamma}_e^i$ 在 $I_j$ 上可写成复合
$\gamma_e^i\circ P_j^i$, 其中 $|(P_j^i)'|=2\pi S^i/(La_j^i)$. 此外,
$P_j^1$ 与 $P_j^2$ 都把 $I_j$ 映入 $[x_{2j},x_{2j+2}]$, 且
\begin{equation}
\int_{I_j}\bigl|(\tilde{\gamma}_e^1-\tilde{\gamma}_e^2)'\bigr|^2
=\int_{I_j}\bigl|(\gamma_e^1\circ P_j^1-\gamma_e^2\circ P_j^2)'\bigr|^2.
\label{eq:5.22}
\end{equation}
最后, 由于速度 $|(P_j^i)'|$ 连续依赖于 $\gamma^i$%
\footnote{速度连续是因为 $a_j^i$ 有正的下界.}, 而 $\gamma_e^i$ 在
$[x_{2j},x_{2j+2}]$ 上 $C^2$ 有界且 $C^1$ 接近, 上式也可变得任意小.
因此, 由于这样的区间至多有 $L$ 个, 在这些区间上的积分也可变得任意小.
\end{proof}

% ================================================================
\section{闭测地线的存在性与宽度}
\label{sec:5.3}

记 $\Omega$ 为连续映射
\begin{equation*}
\sigma:\bS^1\times[0,1]\to M
\end{equation*}
之集, 使得:

\begin{itemize}[leftmargin=1.8em,itemsep=1pt,topsep=2pt]
\item 对每个 $t$, 映射 $\sigma(\cdot,t)$ 属于 $W^{1,2}$.
\item 映射 $t\mapsto\sigma(\cdot,t)$ 是从 $[0,1]$ 到 $W^{1,2}$ 的连续映射.
\item $\sigma$ 把 $\bS^1\times\{0\}$ 与 $\bS^1\times\{1\}$ 映为点.
\end{itemize}

给定映射 $\hat{\sigma}\in\Omega$, 同伦类 $\Omega_{\hat{\sigma}}$ 定义为 $\Omega$
中所有与 $\hat{\sigma}$ 在 $\Omega$ 内同伦的映射之集. 与同伦类
$\Omega_{\hat{\sigma}}$ 相联系的宽度 $W=W(\hat{\sigma})$ 定义为各切片能量的
极大值之下确界. 也就是说, 令
\begin{equation}
W=\inf_{\sigma\in\Omega_{\hat{\sigma}}}\max_{t\in[0,1]}
\mathrm{E}\bigl(\sigma(\cdot,t)\bigr),
\label{eq:5.23}
\end{equation}
其中能量由 $\mathrm{E}(\sigma(\cdot,t))
=\int_{\bS^1}|\partial_x\sigma(x,t)|^2dx$ 给出. 宽度总是非负的, 并且当
$\hat{\sigma}$ 处于非平凡的同伦类中时为正.

\mtfigure[0.82]{fig-5-1}{一个扫掠.}{fig:5.1}

一个特别有趣的例子是 $M$ 为拓扑 $2$-球面、且从 $\bS^2$ 到 $M$ 的诱导映射
的映射度为 $1$ 的情形. 此时宽度为正, 且可由一条非平凡的闭测地线实现.
为看出宽度在非平凡同伦类上为正, 注意若某个切片的极大能量充分小, 则每条
曲线 $\sigma(\cdot,t)$ 都含于 $M$ 的一个凸测地球内. 于是存在一条测地同伦
把 $\sigma$ 与一族点曲线连接起来, 故 $\sigma$ 是同伦平凡的.

本节的主要定理是定理~\ref{thm:5.6}, 它断言收紧后的扫掠中几乎极大的切片
几乎就是测地线; 其证明见 \S 3.4. 该定理的证明以及一列越来越紧的扫掠的
构造都要用到下一小节定义的曲线缩短映射.

\mtpara{宽度关于度量连续}
不难看出宽度关于度量是连续的, 但实现它的极小极大曲线未必如此; 见下面的
例子. 事实上, 把该例子稍加发挥便容易看出, 宽度关于度量一般不会比连续更光滑.

\begin{lemma}
\label{lem:5.4}
宽度在固定闭曲面上的光滑单参数度量族 $\{g_t\}_{t\in[0,1]}$ 上是连续的.
\end{lemma}

\begin{proof}
连续性可由下述事实立即推出:

给定 $\epsilon>0$, 存在 $\delta>0$, 使得若 $t\in[0,1]$ 且 $|s-t|<\delta$,
则
\begin{equation*}
W(g_s)<W(g_t)+\epsilon.
\end{equation*}
\end{proof}

\mtfigure[0.86]{fig-5-2}{Birkhoff 的想法: 把扫掠收紧以得到一条闭测地线.}{fig:5.2}

\mtfigure[0.30]{fig-5-3}{一族单参数``哑铃''度量的三个快照. 实现宽度的
测地线从一个铃铛跳到另一个铃铛. 跳跃发生在中间的图中, 那里测地线不唯一.}%
{fig:5.3}

\mtpara{曲线缩短与扫掠}
下一个结果表明 $\Psi$ 保持扫掠的同伦类.

\begin{lemma}
\label{lem:5.5}
设 $\gamma\in\Omega$ 满足 $\max_t\mathrm{E}(\gamma(\cdot,t))\leq L$. 若
$\gamma_e$ 与 $\tilde{\gamma}_e$ 由对每个 $\gamma(\cdot,t)$ 施行步骤
$(A_1)$ 与 $(A_2)$ 得到, 则 $\gamma$、$\gamma_e$ 与 $\tilde{\gamma}_e$
在 $\Omega$ 中都同伦.
\end{lemma}

\begin{proof}
给定 $x,y\in M$ 且 $\dist_M(x,y)\leq 4\pi$, 如 (C2) 中那样令
$H(x,y):[0,1]\to M$ 为从 $x$ 到 $y$ 的线性映射. 由此
\begin{equation}
F(x,t,s)=H\bigl(\gamma(x,t),\gamma_e(x,t)\bigr)(s)
\label{eq:5.24}
\end{equation}
是一个显式同伦, 且 $F(\cdot,\cdot,0)=\gamma$, $F(\cdot,\cdot,1)=\gamma_e$.

对每个满足 $\Length(\gamma_e(\cdot,t))>0$ 的 $t$, $\gamma_e$ 可写成
$\gamma_e(\cdot,t)=\tilde{\gamma}_e(\cdot,t)\circ P_t$, 其中 $P_t$ 是
$\bS^1$ 的单调重新参数化, 且固定 $x_0=x_{2L}$. 此外, 由 \eqref{eq:5.10}
知 $P_t$ 连续, 而由引理~\ref{lem:5.3} 知 $P_t$ 连续依赖于 $t$. 由于
$x\mapsto(1-s)P_t(x)+sx$ 给出从 $P_t$ 到 $\bS^1$ 上恒等映射的同伦, 我们
断定
\begin{equation}
G(x,t,s)=\tilde{\gamma}_e\bigl((1-s)P_t(x)+sx,t\bigr)
\label{eq:5.25}
\end{equation}
是一个显式同伦, 且 $G(\cdot,\cdot,0)=\gamma_e$,
$G(\cdot,\cdot,1)=\tilde{\gamma}_e$. 注意当
$\Length(\gamma_e(\cdot,t))=0$ 时 $P_t$ 没有定义, 但同伦 $G$ 仍有定义.
\end{proof}

\mtpara{构造扫掠}
选取一列映射 $\hat{\sigma}^j\in\Omega_{\hat{\sigma}}$, 使得
\begin{equation}
\max_{t\in[0,1]}\mathrm{E}\bigl(\hat{\sigma}^j(\cdot,t)\bigr)
<W+\frac1j.
\label{eq:5.26}
\end{equation}
注意 \eqref{eq:5.26} 与 Cauchy--Schwarz 不等式蕴含切片的长度的一致界以及
一致的 $C^{1/2}$ 连续性, 二者都与 $t$ 和 $j$ 无关. 前者立即推出, 后者则由
\begin{equation}
\begin{aligned}
\bigl|\hat{\sigma}^j(x,t)-\hat{\sigma}^j(y,t)\bigr|^2
&\leq\Bigl(\int_x^y\bigl|\partial_s\hat{\sigma}^j(s,t)\bigr|ds\Bigr)^2\\
&\leq|y-x|\int_x^y\bigl|\partial_s\hat{\sigma}^j(s,t)\bigr|^2ds
 \leq|y-x|(W+1)
\end{aligned}
\label{eq:5.27}
\end{equation}
推出.

我们将把 $\hat{\sigma}^j$ 换成一族扫掠 $\sigma^j$, 它们除满足 \eqref{eq:5.26}
之外, 还满足切片 $\sigma^j(\cdot,t)$ 属于 $\Lambda$. 做法是采用与 $\Psi$
的构造中第一步类似的局部线性替换. 也就是说, 切片的一致 $C^{1/2}$ 界
使我们能固定 $\bS^1$ 中点的一个分划 $y_0,\dots,y_N=y_0$, 使得每个区间
$[y_i,y_{i+1}]$ 总被映入 $M$ 中半径至多 $4\pi$ 的球. 其次, 对每个 $t$ 与
每个 $j$, 把 $\hat{\sigma}^j(\cdot,t)|_{[y_i,y_{i+1}]}$ 替换为具有相同端点
的线性映射(测地线), 并把所得映射记为 $\tilde{\sigma}^j(\cdot,t)$. 再
重新参数化 $\tilde{\sigma}^j(\cdot,t)$ 使之具有常速, 得到
$\sigma^j(\cdot,t)$. 容易看出每个 $\sigma^j(\cdot,t)$ 都满足
\eqref{eq:5.26}. 此外, $\sigma^j(\cdot,t)$ 的长度界也给出这些线性映射的
一致 Lipschitz 界; 取 $L$ 为 $N$ 与该 Lipschitz 界中的最大者.

还需验证 $\sigma^j$ 在横向, 即关于 $t$, 连续, 并且在 $\Omega$ 中与
$\hat{\sigma}$ 同伦.\footnote{这些事实 Birkhoff 已经建立; 见 \cite{B1}、
\cite{B2} 以及 \cite{Cr} 第 2 节.} 但这可由引理~\ref{lem:5.3} 与
\ref{lem:5.5} 立即推出.

最后, 对每个 $\sigma^j(\cdot,t)$ 施行替换映射 $\Psi$, 得到一列新的扫掠
$\gamma^j=\Psi(\sigma^j)$. (由引理~\ref{lem:5.3} 与 \ref{lem:5.5}, $\Psi$
连续依赖于 $t$ 并保持同伦类 $\Omega_{\hat{\sigma}}$; 显然 $\Psi$ 固定
$t=0$ 与 $t=1$ 处的常值映射.)

\mtpara{几乎极大蕴含几乎临界}
{\setlength{\emergencystretch}{2em}%
我们证明下述有用的性质; 见下文定理~5.5, 并参看 \cite{CM19}、
\cite{CM27}、\cite{CD} 的命题 3.1、\cite{Pi} 的命题 3.1 以及 \cite{Am3}
的 12.5:\par}

\begin{quote}
收紧后的扫掠中每条长度接近该扫掠中最长曲线长度的曲线, 本身必定接近
一条闭测地线. 特别地, 扫掠中存在接近闭测地线的曲线.
\end{quote}

上述有用的性质在几乎任何用扫掠构造变分问题临界点的工作中都是隐含的, 但
并不总是被记录下来, 因为多数作者只关心一个临界点的存在性.

作为应用, 我们证明这一列扫掠 $\gamma^j$ 是紧的. 确切地说, 我们有
\cite{CM28} 的下述定理:

\begin{theorem}
\label{thm:5.6}
(Colding--Minicozzi, \cite{CM28}) 给定 $W\geq 0$ 与 $\epsilon>0$, 存在
$\delta>0$, 使得若 $j>1/\delta$ 且对某个 $t_0$ 有
\begin{equation}
2\pi\,\mathrm{E}\bigl(\gamma^j(\cdot,t_0)\bigr)
=\Length^2\bigl(\gamma^j(\cdot,t_0)\bigr)>2\pi(W-\delta),
\label{eq:5.28}
\end{equation}
则对这同一个 $j$ 有
$\dist\bigl(\gamma^j(\cdot,t_0),G\bigr)<\epsilon$.
\end{theorem}

{\setlength{\emergencystretch}{2em}%
作为立即可得的推论, 我们得到 $\bS^2$ 上任意度量的非平凡闭测地线的
存在性; 这归功于 Birkhoff. 用调和映射热流给出的另一证明见 \cite{LzWl}.\par}

下一个引理把曲线缩短 $\Psi$ 的性质 (3) 与 (4) 结合起来, 是构造所需扫掠
序列的关键.

\begin{lemma}
\label{lem:5.7}
给定 $W\geq 0$ 与 $\epsilon>0$, 存在 $\delta>0$, 使得若
$\gamma\in\Lambda$ 且
\begin{equation}
2\pi(W-\delta)<\Length^2(\Psi(\gamma))\leq\Length^2(\gamma)
<2\pi(W+\delta),
\label{eq:5.29}
\end{equation}
则 $\dist(\Psi(\gamma),G)<\epsilon$.
\end{lemma}

\begin{proof}
若 $W\leq\epsilon^2/6$, Wirtinger 不等式以 $\delta=\epsilon^2/6$ 给出
引理.

以下设 $W>\epsilon^2/6$. 三角不等式给出
\begin{equation}
\dist(\Psi(\gamma),G)\leq\dist(\Psi(\gamma),\gamma)+\dist(\gamma,G).
\label{eq:5.30}
\end{equation}
由于 $\Psi$ 不会使 $\gamma$ 的长度减少太多, 只要 $\delta$ 充分小,
$\Psi$ 的性质 (4) 就把 $\dist(\gamma,G)$ 界定为 $\epsilon/2$. 同样,
只要 $\delta$ 充分小, $\Psi$ 的性质 (3) 就把
$\dist(\Psi(\gamma),\gamma)$ 界定为 $\epsilon/2$.
\end{proof}

\begin{proof}[定理~\ref{thm:5.6}的证明]
取引理~\ref{lem:5.7} 给出的 $\delta$. 由 \eqref{eq:5.28}、
\eqref{eq:5.26} 并利用 $j>1/\delta$, 得
\begin{equation}
2\pi(W-\delta)<\Length^2\bigl(\gamma^j(\cdot,t_0)\bigr)
\leq\Length^2\bigl(\sigma^j(\cdot,t_0)\bigr)<2\pi(W+\delta).
\label{eq:5.31}
\end{equation}
于是, 由于 $\gamma^j(\cdot,t_0)=\Psi(\sigma^j(\cdot,t_0))$,
引理~\ref{lem:5.7} 给出
$\dist\bigl(\gamma^j(\cdot,t_0),G\bigr)<\epsilon$, 如所欲证.
\end{proof}

\mtpara{参数空间}
我们本可以用任意紧集 $\mathcal{P}$ 作为扫掠中圆周的参数空间, 只要求曲线
在 $\partial\mathcal{P}$ 上为常值(或者 $\partial\mathcal{P}=\varnothing$),
而不必用单位区间 $[0,1]$ 并要求曲线起止于点曲线. 此时记 $\Omega^{\mathcal{P}}$
为连续映射 $\sigma:\bS^1\times\mathcal{P}\to M$ 之集, 使得对每个
$t\in\mathcal{P}$, 曲线 $\sigma(\cdot,t)$ 属于 $W^{1,2}$, 映射
$t\mapsto\sigma(\cdot,t)$ 是从 $\mathcal{P}$ 到 $W^{1,2}$ 的连续映射,
并且最后 $\sigma$ 把 $\partial\mathcal{P}$ 映为点曲线. 给定映射
$\hat{\sigma}\in\Omega^{\mathcal{P}}$, 同伦类
$\Omega_{\hat{\sigma}}^{\mathcal{P}}\subset\Omega^{\mathcal{P}}$ 定义为
$\Omega^{\mathcal{P}}$ 中所有与 $\hat{\sigma}$ 在
$\Omega^{\mathcal{P}}$ 内同伦的映射之集. 最后, 宽度
$W=W(\hat{\sigma})$ 为
\begin{equation}
W=\inf_{\sigma\in\Omega_{\hat{\sigma}}^{\mathcal{P}}}
\max_{t\in\mathcal{P}}\mathrm{E}\bigl(\sigma(\cdot,t)\bigr).
\label{eq:5.32}
\end{equation}
定理~\ref{thm:5.6} 对这些一般的参数空间同样成立; 证明几乎相同, 只需作
平凡的改动.

% ================================================================
\section{调和替换}
\label{sec:5.4}

下面转向从 $\bS^2$ 出发的映射. 粗略的想法与曲线情形相同: 我们定义一个能量
不增的``离散梯度流'', 并证明它在远离调和映射处使能量严格减少. 这个能量
不增的映射同样由局部替换来构造, 只不过这次用的是小球上能量很小的调和
(极小化能量)映射作替换. 这通常称为\newterm{harmonic-replacement}{调和替换}%
{harmonic replacement}, 它使人想起 H. A. Schwarz 求解调和函数
Dirichlet 问题的交替法.

\mtpara{调和替换的性质}
本小节总结调和替换的关键性质. 这些性质在 \cite{CM27} 中证明, 其证明请参看
原文.

使调和替换有用的第一个关键性质是: 能量泛函在小能量映射上是严格凸的. 这是
关于调和映射的一个凸性定理. 其确切陈述由 Colding--Minicozzi 的下述结果
给出, 见 \cite{CM27}:

\begin{theorem}
\label{thm:5.8}
(Colding--Minicozzi, \cite{CM27}) 存在常数 $\epsilon_1>0$(依赖于 $M$),
使得若 $u$ 与 $v$ 是从 $B_1\subset\R^2$ 到 $M$ 的 $W^{1,2}$ 映射, $u$
与 $v$ 在 $\partial B_1$ 上相等, 且 $v$ 弱调和而能量至多 $\epsilon_1$, 则
\begin{equation}
\int_{B_1}|\nabla u|^2-\int_{B_1}|\nabla v|^2
\geq\frac12\int_{B_1}|\nabla v-\nabla u|^2.
\label{eq:5.33}
\end{equation}
\end{theorem}

定理~\ref{thm:5.8} 的一个直接推论是小能量映射的 Dirichlet 问题解的
唯一性(以及任何这样的调和映射都极小化能量). 该定理依赖于 H\'elein
\cite{Hf1}、\cite{Hf2} 关于弱调和映射的估计, 以及可追溯到 Wente
\cite{W} 工作的补偿紧性技巧.

\begin{corollary}
\label{cor:5.9}
(Colding--Minicozzi, \cite{CM27}) 设 $\epsilon_1>0$ 如定理~\ref{thm:5.8}
中所示. 若 $u_1$ 与 $u_2$ 是从 $B_1\subset\R^2$ 到 $M$ 的 $W^{1,2}$
弱调和映射, 能量都不超过 $\epsilon_1$, 且它们在 $\partial B_1$ 上相等,
则 $u_1=u_2$.
\end{corollary}

定理~\ref{thm:5.8} 的第二个推论是: 若限制在小能量映射上, 则调和替换作为
从 $C^0(\overline{B_1})\cap W^{1,2}(B_1)$ 到自身的映射是连续的. (在
$C^0(\overline{B_1})\cap W^{1,2}(B_1)$ 上的范数取为上确界范数与
$W^{1,2}$ 范数之和.)

\begin{corollary}
\label{cor:5.10}
设 $\epsilon_1>0$ 如定理~\ref{thm:5.8} 中所示, 并令
\begin{equation}
\mathcal{M}=\Bigl\{u\in C^0(\overline{B_1},M)\cap W^{1,2}(B_1,M)
\mid \mathrm{E}(u)\leq\epsilon_1\Bigr\}.
\label{eq:5.34}
\end{equation}
给定 $u\in\mathcal{M}$, 存在唯一的极小化能量的映射 $w$, 它在
$\partial B_1$ 上等于 $u$, 且 $w\in\mathcal{M}$. 此外, 存在依赖于 $M$
的 $C$, 使得若 $u_1,u_2\in\mathcal{M}$ 对应的极小化能量映射为 $w_1,w_2$,
并令 $E=\mathrm{E}(u_1)+\mathrm{E}(u_2)$, 则
\begin{equation}
\begin{aligned}
\bigl|\mathrm{E}(w_1)-\mathrm{E}(w_2)\bigr|
&\leq C\|u_1-u_2\|_{C^0(\overline{B_1})}E\\
&\quad +C\|\nabla u_1-\nabla u_2\|_{L^2(B_1)}E^{1/2}.
\end{aligned}
\label{eq:5.35}
\end{equation}
最后, 从 $u$ 到 $w$ 的映射作为从
$C^0(\overline{B_1})\cap W^{1,2}(B_1)$ 到自身的映射是连续的.
\end{corollary}

为叙述下一个引理, 引进一些记号是方便的. 确切地说, 给定从 $\bS^2$ 到 $M$
的 $C^0\cap W^{1,2}$ 映射 $u$ 以及 $\bS^2$ 中一族有限个互不相交的闭球
$\mathcal{B}$, 使得 $u$ 在 $\bigcup_{\mathcal{B}}B$ 上的能量至多
$\epsilon_1/3$, 记 $H(u,\mathcal{B}):\bS^2\to M$ 为这样的映射: 它在
$\bS^2\setminus\bigcup_{\mathcal{B}}B$ 上与 $u$ 重合, 而在
$\bigcup_{\mathcal{B}}B$ 上等于从 $\bigcup_{\mathcal{B}}B$ 到 $M$ 的、
在 $\bigcup_{\mathcal{B}}\partial B$ 上与 $u$ 相等的极小化能量映射. 为使
记号简洁, 我们约定
\begin{equation}
H(u,\mathcal{B}_1,\mathcal{B}_2)=H\bigl(H(u,\mathcal{B}_1),\mathcal{B}_2\bigr).
\label{eq:5.36}
\end{equation}
最后, 若 $\alpha\in(0,1]$, 则 $\alpha\mathcal{B}$ 表示由同心球构成、但半径
收缩 $\alpha$ 倍的族.

一般而言, $H(u,\mathcal{B}_1,\mathcal{B}_2)$ 与
$H(u,\mathcal{B}_2,\mathcal{B}_1)$ 并不相同. 这在定理~\ref{thm:5.23} 的
证明中是重要的, 那里在 $\frac12\mathcal{B}_1$ 或 $\frac12\mathcal{B}_2$
上作调和替换都会使 $u$ 的能量下降一个确定的量. 下一个引理(见
\eqref{eq:5.37})表明, 无论作替换的次序如何, 能量都下降一个确定的量.
第二个不等式则用对 $u$ 作调和替换时可能的能量下降, 来界定对
$H(u,\mathcal{B}_1)$ 作调和替换时可能的能量下降.

\begin{lemma}
\label{lem:5.11}
存在常数 $\kappa>0$(依赖于 $M$), 使得若 $u:\bS^2\to M$ 属于
$C^0\cap W^{1,2}$, 且 $\mathcal{B}_1$、$\mathcal{B}_2$ 各是 $\bS^2$ 中
一族有限个互不相交的闭球, 使 $u$ 在每个 $\bigcup_{\mathcal{B}_i}B$ 上的
能量至多 $\epsilon_1/3$, 则
\begin{equation}
\mathrm{E}(u)-\mathrm{E}\bigl[H(u,\mathcal{B}_1,\mathcal{B}_2)\bigr]
\geq\kappa\Bigl(\mathrm{E}(u)-\mathrm{E}\bigl[H(u,\tfrac12\mathcal{B}_2)\bigr]
\Bigr)^2.
\label{eq:5.37}
\end{equation}
此外, 对任意 $\mu\in[1/8,1/2]$ 有
\begin{equation}
\begin{aligned}
&\frac{\bigl(\mathrm{E}(u)-\mathrm{E}[H(u,\mathcal{B}_1)]\bigr)^{1/2}}{\kappa}
+\mathrm{E}(u)-\mathrm{E}\bigl[H(u,2\mu\mathcal{B}_2)\bigr]\\
&\qquad\geq\mathrm{E}\bigl[H(u,\mathcal{B}_1)\bigr]
-\mathrm{E}\bigl[H(u,\mathcal{B}_1,\mu\mathcal{B}_2)\bigr].
\end{aligned}
\label{eq:5.38}
\end{equation}
\end{lemma}

\mtpara{几乎调和映射的紧性}
我们的``几乎调和''概念依赖于从 $\bS^2$ 到 $M$ 的调和映射的两条重要性质.
第一, 从 $\bS^2$ 出发的调和映射是共形的, 因而能量与面积相等; 见下文 (A).
第二, 从曲面出发的任何调和映射限制在能量充分小的球上时都极小化能量;
见下文 (B).

在下面的命题中, $\epsilon_{SU}>0$(依赖于 $M$)是\thmref{thm:4.22} 中的小能量常数,
于是能量至多 $\epsilon_{SU}$ 的调和映射具有内估计. 特别地, 从 $\bS^2$ 到
$M$ 的任何非常值调和映射的能量都大于 $\epsilon_{SU}$.

\begin{proposition}
\label{prop:5.12}
设 $\epsilon_0,E_0>0$ 是常数, $\epsilon_{SU}>\epsilon_0$, 而
$u^j:\bS^2\to M$ 是一列 $C^0\cap W^{1,2}$ 映射, 满足
$E_0\geq\mathrm{E}(u^j)$ 以及:

\begin{enumerate}[label=(\Alph*),leftmargin=2.2em]
\item $\Area(u^j)>\mathrm{E}(u^j)-1/j$.
\item 对 $\bS^2$ 中任意一族有限个互不相交的闭球 $\mathcal{B}$, 只要
      $\int_{\bigcup_{\mathcal{B}}B}|\nabla u^j|^2<\epsilon_0$, 就存在
      极小化能量的映射 $v:\bigcup_{\mathcal{B}}\frac18B\to M$, 它在
      $\bigcup_{\mathcal{B}}\frac18\partial B$ 上等于 $u^j$, 且满足
\begin{equation*}
\int_{\bigcup_{\mathcal{B}}\frac18B}\bigl|\nabla u^j-\nabla v\bigr|^2
\leq 1/j.
\end{equation*}
\end{enumerate}

若 (A) 与 (B) 成立, 则 $u^j$ 的某个子列变分流形收敛到一族调和映射
$v^0,\dots,v^m:\bS^2\to M$.
\end{proposition}

变分流形收敛在\textbf{第 3 章}中定义.

命题~\ref{prop:5.12} 的一个直接推论是有界能量调和映射序列的紧性定理.
这由 Jost 在 \cite{Jo} 的引理 4.3.1 中证明. 事实上, Parker 在更强的拓扑
下证明了有界能量调和映射的紧性, 即除 $W^{1,2}$ 收敛之外还有 $C^0$ 收敛;
见 \cite{Pk} 的定理 2.2. 因此, 对上述意义下越来越接近调和的序列也有类似
的紧性, 也许并不令人意外. 然而, 记住 Parker 构造过一些映射序列, 其
Laplace 算子在 $L^1$ 中趋于零却没有收敛子列(见 \cite{Pk} 的命题 4.2),
这是有益的.

最后指出, 命题~\ref{prop:5.12} 可以看作 Palais--Smale 条件 (C) 的离散
版本. 也就是说, 若一列映射由调和替换所能带来的最大能量下降趋于零, 则
某个子列收敛到一族调和映射.

命题~\ref{prop:5.12} 的证明将遵循 Parker 与 Wolfson 在 \cite{PkW} 中发展、
并被 Parker 在 \cite{Pk} 中用来证明有界能量调和映射紧性的一般结构. 主要
困难在于排除极限中的能量损失(见气泡收敛定义中的 (B4); 这称为
\newterm{energy-identity}{能量等式}{energy identity}). 处理这个问题的
粗略想法是: 能量损失只发生在很小的圆环上, 那里映射``几乎''调和, 而圆环
的内外半径之比极大. 我们将用命题~\ref{prop:5.18} 证明映射在这样一个圆环
上必定``远离''共形, 于是条件 (A) 使我们能够排除能量损失. 这里``远离''
共形是指映射的 $\theta$-能量远小于径向能量. 为使之精确, 宜于把 $\R^2$
中的圆环 $B_{e^{r_2}}\setminus B_{e^{r_1}}$ 换为共形等价的柱面
$[r_1,r_2]\times\bS^1$. 带有平坦乘积度量、坐标为 $t$ 与 $\theta$ 的
(非紧)柱面 $\R\times\bS^1$ 将记为 $\mathcal{C}$. 对 $r_1<r_2$, 记
$\mathcal{C}_{r_1,r_2}\subset\mathcal{C}$ 为乘积
$[r_1,r_2]\times\bS^1$.

\mtpara{柱面上的调和映射}
本小节的主要结果是: 长柱面上能量很小的调和映射几乎是径向的. 这意味着
一列能量有正下界的这种映射一致地远离共形, 因而不能对命题~\ref{prop:5.12}
中的 (A) 成立. 这个结果将用来在命题~\ref{prop:5.18} 中对``几乎调和''
映射证明类似结论, 并最终用于证明能量不会损失.

\begin{proposition}
\label{prop:5.13}
给定 $\delta>0$, 存在依赖于 $\delta$(以及 $M$)的 $\epsilon_2>0$ 与
$\ell\geq 1$, 使得若 $u$ 是从平坦柱面
$\mathcal{C}_{-3\ell,3\ell}=[-3\ell,3\ell]\times\bS^1$ 到 $M$ 的(非常值)
$C^3$ 调和映射且 $\mathrm{E}(u)\leq\epsilon_2$, 则
\begin{equation}
\int_{\mathcal{C}_{-\ell,\ell}}|u_\theta|^2
<\delta\int_{\mathcal{C}_{-2\ell,2\ell}}|\nabla u|^2.
\label{eq:5.39}
\end{equation}
\end{proposition}

为证明这个命题, 我们证明一个微分不等式, 它导致柱面各水平集上调和映射的
$\theta$-能量指数增长. 一旦有了这一点, 命题便随之推出. 也就是说, 若柱面
``中部''的 $\theta$-能量是双倍柱面上总能量的一个确定比例, 则指数增长
将迫使柱面边界附近的 $\theta$-能量过大.

下述标准引理正是导致 $\theta$-能量(通过下文引理~\ref{lem:5.15})指数
增长的微分不等式.

\begin{lemma}
\label{lem:5.14}
对从 $\mathcal{C}_{r_1,r_2}\subset\mathcal{C}$ 到 $M\subset\R^N$ 的
$C^3$ 调和映射 $u$,
\begin{equation}
\partial_t^2\int_t|u_\theta|^2
\geq\frac32\int_t|u_\theta|^2-2\sup_M|A|^2\int_t|\nabla u|^4.
\label{eq:5.40}
\end{equation}
\end{lemma}

\begin{proof}
对 $\int_t|u_\theta|^2$ 求导并在 $\theta$ 方向分部积分, 得
\begin{equation}
\begin{aligned}
\frac12\partial_t^2\int_t|u_\theta|^2
&=\int_t|u_{t\theta}|^2+\int_t\langle u_\theta,u_{tt\theta}\rangle
 =\int_t|u_{t\theta}|^2-\int_t\langle u_{\theta\theta},u_{tt}\rangle\\
&=\int_t|u_{t\theta}|^2-\int_t\langle u_{\theta\theta},
 (\Delta u-u_{\theta\theta})\rangle\\
&\geq\int_t|u_{t\theta}|^2+\int_t|u_{\theta\theta}|^2
 -\sup_M|A|\int_t|u_{\theta\theta}|\,|\nabla u|^2,
\end{aligned}
\label{eq:5.41}
\end{equation}
其中最后一个不等式用了调和映射方程 \eqref{eq:4.34} 给出的
$|\Delta u|\leq|\nabla u|^2\sup_M|A|$. 应用吸收不等式
$2ab\leq a^2/2+2b^2$, 并注意 $\int_tu_\theta=0$ 从而由 Wirtinger
不等式有 $\int_t|u_\theta|^2\leq\int_t|u_{\theta\theta}|^2$, 引理便得证.
\end{proof}

我们还需要一个简单的常微分方程比较引理:

\begin{lemma}
\label{lem:5.15}
设 $f$ 是 $[-2\ell,2\ell]\subset\R$ 上的非负 $C^2$ 函数, 满足
\begin{equation}
f''\geq f-a,
\label{eq:5.42}
\end{equation}
其中 $a>0$ 为常数. 若 $\max_{[-\ell,\ell]}f\geq 2a$, 则
\begin{equation}
\int_{-2\ell}^{2\ell}f\geq 2\sqrt2\,a\sinh\bigl(\ell/\sqrt2\bigr).
\label{eq:5.43}
\end{equation}
\end{lemma}

\begin{proof}
取 $x_0\in[-\ell,\ell]$ 使 $f$ 在 $[-\ell,\ell]$ 上达到最大值. 由于引理
在反射 $x\mapsto-x$ 下不变, 可设 $x_0\geq 0$. 若 $x_0$ 是内点, 则
$f'(x_0)=0$; 否则若 $x_0=\ell$, 则 $f'(x_0)\geq 0$. 无论哪种情形都有
$f'(x_0)\geq 0$. 由于 $f(x_0)\geq 2a$, \eqref{eq:5.42} 给出
$f''(x_0)\geq a>0$, 因而 $f'$ 在 $x_0$ 处严格递增.

我们断言: 对 $(x_0,2\ell]$ 中所有 $x$ 都有 $f'(x)>0$. 否则将存在第一个
点 $y>x_0$ 使 $f'(y)=0$. 由此 $f'\geq 0$ 于 $[x_0,y]$, 于是
$f\geq f(x_0)\geq 2a$ 于 $[x_0,y]$, 从而 $f''\geq a>0$ 于 $[x_0,y]$,
这与 $f'(y)\leq f'(x_0)$ 矛盾.

由该断言, $f$ 在 $[x_0,2\ell]$ 上单调递增, 于是 \eqref{eq:5.42} 给出
\begin{equation}
f''\geq\frac12f\qquad\text{于 }[x_0,2\ell].
\label{eq:5.44}
\end{equation}
利用 $f'(x_0)\geq 0$ 与 \eqref{eq:5.44} 作标准的 Riccati 比较论证
(例如见 \cite{CDM} 的推论 A.9), 对 $t\in[0,2\ell-x_0]$ 得到
\begin{equation}
f(x_0+t)\geq f(x_0)\cosh\bigl(t/\sqrt2\bigr)
\geq 2a\cosh\bigl(t/\sqrt2\bigr).
\label{eq:5.45}
\end{equation}
最后, 把 \eqref{eq:5.45} 在 $[0,\ell]$ 上积分即得 \eqref{eq:5.43}.
\end{proof}

\begin{proof}[命题~\ref{prop:5.13}的证明]
由于我们将取 $\ell\geq 1$ 与 $\epsilon_2<\epsilon_{SU}$, 调和映射的小
能量内估计(见\thmref{thm:4.22})蕴含
\begin{equation}
\sup_{\mathcal{C}_{-2\ell,2\ell}}|\nabla u|^2
\leq C_{SU}\int_{\mathcal{C}_{-3\ell,3\ell}}|\nabla u|^2
\leq C_{SU}\epsilon_2.
\label{eq:5.46}
\end{equation}
令 $f(t)=\int_t|u_\theta|^2$. 由引理~\ref{lem:5.14} 推出
\begin{equation}
\begin{aligned}
f''(t)
&\geq\frac32f(t)-2\sup_M|A|^2C_{SU}\epsilon_2
 \int_t\bigl(|u_\theta|^2+|u_t|^2\bigr)\\
&\geq f(t)-C\epsilon_2\int_t\bigl(|u_t|^2-|u_\theta|^2\bigr),
\end{aligned}
\label{eq:5.47}
\end{equation}
其中 $C=2C_{SU}\sup_M|A|^2$, 且第二个不等式用了假设
$C\epsilon_2\leq 1/4$.

我们将用到 $\int_t(|u_t|^2-|u_\theta|^2)$ 关于 $t$ 为常数. 为看出这一点,
求导得
\begin{equation}
\begin{aligned}
\frac12\partial_t\int_t\bigl(|u_t|^2-|u_\theta|^2\bigr)
&=\int_t\bigl(\langle u_t,u_{tt}\rangle-\langle u_\theta,u_{t\theta}\rangle\bigr)\\
&=\int_t\bigl\langle u_t,(u_{tt}+u_{\theta\theta})\bigr\rangle=0,
\end{aligned}
\label{eq:5.48}
\end{equation}
其中第二个等式用了在 $\theta$ 方向的分部积分, 最后一个等式用了
$u_{tt}+u_{\theta\theta}=\Delta u$ 垂直于 $M$ 而 $u_t$ 与 $M$ 相切.
把这个常数界定为
\begin{equation}
\int_t\bigl(|u_t|^2-|u_\theta|^2\bigr)
=\frac1{4\ell}\int_{\mathcal{C}_{-2\ell,2\ell}}
\bigl(|u_t|^2-|u_\theta|^2\bigr)
\leq\frac1{4\ell}\int_{\mathcal{C}_{-2\ell,2\ell}}|\nabla u|^2.
\label{eq:5.49}
\end{equation}
由 \eqref{eq:5.47} 与 \eqref{eq:5.49}, 对
$a=\frac{C\epsilon_2}{4\ell}\int_{\mathcal{C}_{-2\ell,2\ell}}|\nabla u|^2$
应用引理~\ref{lem:5.15}, 蕴含或者
\begin{equation}
\max_{[-\ell,\ell]}f<2\,\frac{C\epsilon_2}{4\ell}
\int_{\mathcal{C}_{-2\ell,2\ell}}|\nabla u|^2,
\label{eq:5.50}
\end{equation}
或者
\begin{equation}
\int_{\mathcal{C}_{-2\ell,2\ell}}|u_\theta|^2
=\int_{-2\ell}^{2\ell}f(t)\,dt
\geq 2\sqrt2\,C\epsilon_2\frac{\sinh(\ell/\sqrt2)}{4\ell}
\int_{\mathcal{C}_{-2\ell,2\ell}}|\nabla u|^2.
\label{eq:5.51}
\end{equation}
只要 $\ell$ 充分大, 使得
\begin{equation}
2\sqrt2\,C\epsilon_2\frac{\sinh(\ell/\sqrt2)}{4\ell}>1,
\label{eq:5.52}
\end{equation}
第二种可能性便不会发生.

用 \eqref{eq:5.50} 作为 $f$ 在 $[-\ell,\ell]$ 上的上界来估计 $f$ 的积分,
给出
\begin{equation}
\int_{\mathcal{C}_{-\ell,\ell}}|u_\theta|^2
\leq 2\ell\max_{[-\ell,\ell]}f
<C\epsilon_2\int_{\mathcal{C}_{-2\ell,2\ell}}|\nabla u|^2.
\label{eq:5.53}
\end{equation}
于是, 取 $\epsilon_2>0$ 使 $C\epsilon_2<\min\{1/4,\delta\}$, 再取 $\ell$
使 \eqref{eq:5.52} 成立, 命题便得证.
\end{proof}

\mtpara{几乎调和映射的弱紧性}
我们将需要关于 $W^{1,2}(\bS^2,M)$ 中一列能量一致有界、且局部被调和映射
很好地逼近的映射的紧性定理. 在精确陈述之前, 回顾调和映射的情形是有益的.
因此设 $u^j:\bS^2\to M$ 是一列调和映射, $\mathrm{E}(u^j)\leq E_0$, 其中
$E_0$ 固定. 取子列后可设测度 $|\nabla u^j|^2dx$ 收敛, 且存在一个有限点集
$\mathcal{S}$, 能量集中于此, 使得
\begin{equation}
\text{若 }x\in\mathcal{S},\ \text{则 }
\inf_{r>0}\Bigl[\lim_{j\to\infty}\int_{B_r(x)}|\nabla u^j|^2\Bigr]
\geq\epsilon_{SU},
\label{eq:5.54}
\end{equation}
\begin{equation}
\text{若 }x\notin\mathcal{S},\ \text{则 }
\inf_{r>0}\Bigl[\lim_{j\to\infty}\int_{B_r(x)}|\nabla u^j|^2\Bigr]
<\epsilon_{SU}.
\label{eq:5.55}
\end{equation}
常数 $\epsilon_{SU}>0$ 来自\thmref{thm:4.22}, 于是 \eqref{eq:5.55} 蕴含 $u^j$ 在
$x$ 的某个邻域内具有一致的 $C^{2,\alpha}$ 估计. 因此, Arzel\`a--Ascoli
定理与对角论证给出 $u^j$ 的进一步子列, 它在 $\bS^2\setminus\mathcal{S}$
的每个紧子集上 $C^2$ 收敛到一张调和映射. 我们将需要这一事实的更一般版本,
其中 $u^j:\bS^2\to M$ 是一列 $\mathrm{E}(u^j)\leq E_0$ 的 $W^{1,2}$ 映射,
且在下述意义下 $\epsilon_0$-几乎调和:

$(B_0)$ 若 $B\subset\bS^2$ 是任一满足
$\int_B|\nabla u^j|^2<\epsilon_0$ 的球, 则存在极小化能量的映射
$v:\frac18B\to M$, 它在 $\partial\frac18B$ 上与 $u^j$ 取相同的边界值, 且
\begin{equation*}
\int_{\frac18B}\bigl|\nabla u^j-\nabla v\bigr|^2\leq 1/j.
\end{equation*}

\begin{lemma}
\label{lem:5.16}
设 $\epsilon_0>0$ 小于 $\epsilon_{SU}$. 若 $u^j:\bS^2\to M$ 是一列满足
$(B_0)$ 且 $\mathrm{E}(u^j)\leq E_0$ 的 $W^{1,2}$ 映射, 则存在有限点集
$\{x_1,\dots,x_k\}$、仍记为 $u^j$ 的子列以及调和映射 $u:\bS^2\to M$,
使得 $u^j$ 在 $W^{1,2}$ 中弱收敛到 $u$, 并且若
$K\subset\bS^2\setminus\{x_1,\dots,x_k\}$ 是紧集, 则 $u^j$ 在
$W^{1,2}(K)$ 中收敛到 $u$. 此外, 测度 $|\nabla u^j|^2dx$ 收敛到测度
$\nu$, 满足 $\epsilon_0\leq\nu(x_i)$ 与 $\nu(\bS^2)\leq E_0$.
\end{lemma}

\begin{proof}
取子列后可设:

\begin{itemize}[leftmargin=1.8em,itemsep=1pt,topsep=2pt]
\item $u^j$ 在 $W^{1,2}$ 中弱收敛到 $W^{1,2}$ 映射 $u:\bS^2\to M$.
\item 测度 $|\nabla u^j|^2dx$ 收敛到极限测度 $\nu$, 且
      $\nu(\bS^2)\leq E_0$.
\end{itemize}

由此, 满足
\begin{equation}
\lim_{r\to 0}\nu\bigl(B_r(x_j)\bigr)\geq\epsilon_0
\label{eq:5.56}
\end{equation}
的点 $x_1,\dots,x_k$ 至多有 $E_0/\epsilon_0$ 个.

下面证明在远离这些 $x_i$ 处收敛是 $W^{1,2}$ 中的强收敛, 且 $u$ 是调和的.
为此考虑点 $x\notin\{x_1,\dots,x_k\}$. 按定义, 存在 $r_x>0$ 与 $J_x$,
使得当 $j\geq J_x$ 时 $\int_{B_{r_x}(x)}|\nabla u^j|^2<\epsilon_0$. 特别地,
$(B_0)$ 可以应用, 于是得到极小化能量的映射
$v_x^j:\frac18B_{r_x}(x)\to M$, 它在
$\partial\frac18B_{r_x}(x)$ 上与 $u^j$ 相等, 且满足
\begin{equation}
\int_{\frac18B_{r_x}(x)}\bigl|\nabla v_x^j-\nabla u^j\bigr|^2\leq 1/j.
\label{eq:5.57}
\end{equation}
(这里 $\frac18B_{r_x}(x)$ 是 $\bS^2$ 中以 $x$ 为心的球, 使得把 $x$ 送到
$0\in\R^2$ 的球极投影 $\Pi_x$ 把 $\frac18B_{r_x}(x)$ 与 $B_{r_x}(x)$
映为以 $0$ 为心、半径相差 $8$ 倍的球.) 由于
$\mathrm{E}(v_x^j)\leq\epsilon_0\leq\epsilon_{SU}$, 由\thmref{thm:4.22} 知
$v_x^j$ 的某个子列在 $W^{1,2}(\frac19B_{r_x}(x))$ 中强收敛到调和映射
$v_x:\frac19B_{r_x}(x)\to M$. 把它与三角不等式及 \eqref{eq:5.57}
结合起来, 得
\begin{equation}
\begin{aligned}
\int_{\frac19B_{r_x}(x)}\bigl|\nabla u^j-\nabla v_x\bigr|^2
&\leq 2\int_{\frac19B_{r_x}(x)}\bigl|\nabla u^j-\nabla v_x^j\bigr|^2\\
&\quad +2\int_{\frac19B_{r_x}(x)}\bigl|\nabla v_x^j-\nabla v_x\bigr|^2\to 0.
\end{aligned}
\label{eq:5.58}
\end{equation}
类似地, 由这一收敛、三角不等式、\eqref{eq:5.57} 以及 Dirichlet--Poincar\'e
不等式(\cite{Ev} 第 265 页定理 3; 它适用是因为 $v_x^j$ 在
$\partial\frac18B_{r_x}(x)$ 上等于 $u^j$)给出
\begin{equation}
\begin{aligned}
\int_{\frac19B_{r_x}(x)}\bigl|u^j-v_x\bigr|^2
&\leq 2\int_{\frac18B_{r_x}(x)}\bigl|u^j-v_x^j\bigr|^2\\
&\quad +2\int_{\frac19B_{r_x}(x)}\bigl|v_x^j-v_x\bigr|^2\to 0.
\end{aligned}
\label{eq:5.59}
\end{equation}
结合 \eqref{eq:5.58} 与 \eqref{eq:5.59} 可见 $u^j$ 在
$W^{1,2}(\frac19B_{r_x}(x))$ 中强收敛到 $v_x$. 特别地,
$u|_{\frac19B_{r_x}(x)}=v_x$. 于是 $u$ 在
$\bS^2\setminus\{x_1,\dots,x_k\}$ 上调和. 此外, 由于任一紧集
$K\subset\bS^2\setminus\{x_1,\dots,x_k\}$ 都可由有限个这样的九分之一球
覆盖, 故 $u^j$ 在 $W^{1,2}(K)$ 中强收敛到 $u$.

最后, 由于 $u$ 能量有限, 它在每个 $x_i$ 处必有可去奇点, 因而可延拓为整个
$\bS^2$ 上的调和映射(见 \cite{SaUh} 的定理 3.6).
\end{proof}

\mtpara{柱面上的几乎调和映射}
本小节的主要结果, 即下文命题~\ref{prop:5.18}, 把命题~\ref{prop:5.13} 从
调和映射推广到``几乎调和''映射. 这里``几乎调和''由下文定义~\ref{def:5.17}
精确化, 粗略地说, 它是指在一定的球上作调和替换不会使能量减少很多.

\begin{definition}
\label{def:5.17}
给定 $\nu>0$ 与柱面 $\mathcal{C}_{r_1,r_2}$, 我们称
$W^{1,2}(\mathcal{C}_{r_1,r_2},M)$ 映射 $u$ 是 \textbf{$\nu$-几乎调和}
的, 若对共形等价的圆环
$B_{e^{r_2}}\setminus B_{e^{r_1}}\subset\R^2$ 中任意一族有限个互不
相交的闭球 $\mathcal{B}$, 都存在极小化能量的映射
$v:\bigcup_{\mathcal{B}}\frac18B\to M$, 它在
$\bigcup_{\mathcal{B}}\frac18\partial B$ 上等于 $u$, 且满足
\begin{equation}
\int_{\bigcup_{B}\frac18B}|\nabla u-\nabla v|^2
\leq\frac{\nu}{2}\int_{\mathcal{C}_{r_1,r_2}}|\nabla u|^2.
\label{eq:5.60}
\end{equation}
\end{definition}

这里我们稍微滥用了记号, 因为我们的集合总被看作柱面的子集; 也就是说, 我们把
圆环中的欧氏球与其在共形映射下的像等同起来.

在本小节及其后两小节中, 给定 $\delta>0$, 常数 $\ell\geq 1$ 与
$\epsilon_2>0$ 由命题~\ref{prop:5.13} 给出; 它们只依赖于 $M$ 与 $\delta$.

\begin{proposition}
\label{prop:5.18}
给定 $\delta>0$, 存在 $\nu>0$(依赖于 $\delta$ 与 $M$), 使得若 $m$ 是
正整数, $u$ 是从 $\mathcal{C}_{-(m+3)\ell,3\ell}$ 到 $M$ 的
$\nu$-几乎调和映射且 $\mathrm{E}(u)\leq\epsilon_2$, 则
\begin{equation}
\int_{\mathcal{C}_{-m\ell,0}}|u_\theta|^2
\leq 7\delta\int_{\mathcal{C}_{-(m+3)\ell,3\ell}}|\nabla u|^2.
\label{eq:5.61}
\end{equation}
\end{proposition}

我们将用紧性论证把命题~\ref{prop:5.18} 化归到调和映射的情形, 然后诉诸
命题~\ref{prop:5.13} 来证明它. 一个关键困难是命题~\ref{prop:5.18} 中对
柱面的长度(即对 $m$)没有上界, 因此不能直接应用紧性论证. 这可以通过把
柱面分成固定大小的子柱面再作覆盖论证来处理.

\mtpara{紧性论证}
下一个引理把命题~\ref{prop:5.13} 从 $\mathcal{C}_{-3\ell,3\ell}$ 上的
调和映射推广到几乎调和映射. 与命题~\ref{prop:5.18} 的主要区别在于引
理~\ref{lem:5.19} 中柱面的大小是固定的.

\begin{lemma}
\label{lem:5.19}
给定 $\delta>0$, 存在 $\mu>0$(依赖于 $\delta$ 与 $M$), 使得若 $u$ 是从
$\mathcal{C}_{-3\ell,3\ell}$ 到 $M$ 的 $\mu$-几乎调和映射且
$\mathrm{E}(u)\leq\epsilon_2$, 则
\begin{equation}
\int_{\mathcal{C}_{-\ell,\ell}}|u_\theta|^2
\leq\delta\int_{\mathcal{C}_{-3\ell,3\ell}}|\nabla u|^2.
\label{eq:5.62}
\end{equation}
\end{lemma}

\begin{proof}
我们用反证法, 设存在一列 $1/j$-几乎调和映射
$u^j:\mathcal{C}_{-3\ell,3\ell}\to M$, 满足
$\mathrm{E}(u^j)\leq\epsilon_2$ 且
\begin{equation}
\int_{\mathcal{C}_{-\ell,\ell}}|u_\theta^j|^2
>\delta\int_{\mathcal{C}_{-3\ell,3\ell}}|\nabla u^j|^2.
\label{eq:5.63}
\end{equation}
我们将证明 $u^j$ 的某个子列收敛到一张非常值的调和映射, 与
命题~\ref{prop:5.13} 矛盾. 按 $\mathrm{E}(u^j)$ 是否趋于 $0$ 分两种情形
讨论.

先设 $\limsup_{j\to\infty}\mathrm{E}(u^j)>0$. 能量上界与
$1/j$-几乎调和的假设(以及 $M$ 的紧性)使我们能像引理~\ref{lem:5.16} 中
那样论证, 得到子列在 $\mathcal{C}_{-3\ell,3\ell}$ 的紧子集上 $W^{1,2}$
收敛到一张\emph{非常值}调和映射
$\tilde{u}:\mathcal{C}_{-3\ell,3\ell}\to M$. 此外, 利用
$\mathcal{C}_{-\ell,\ell}$ 上的 $W^{1,2}$ 收敛以及能量的下半连续性,
\eqref{eq:5.63} 蕴含
\begin{equation}
\int_{\mathcal{C}_{-\ell,\ell}}|\tilde{u}_\theta|^2
\geq\delta\int_{\mathcal{C}_{-3\ell,3\ell}}|\nabla\tilde{u}|^2.
\label{eq:5.64}
\end{equation}
这与命题~\ref{prop:5.13} 矛盾.

现设 $\mathrm{E}(u^j)\to 0$. 把 $u^j$ 换为
\begin{equation}
v^j=\frac{u^j-u^j(0)}{\bigl(\mathrm{E}(u^j)\bigr)^{1/2}},
\label{eq:5.65}
\end{equation}
便得到一列映入
\begin{equation}
M_j=\frac{M-u^j(0)}{\bigl(\mathrm{E}(u^j)\bigr)^{1/2}}
\label{eq:5.66}
\end{equation}
的映射, 且 $\mathrm{E}(v^j)=1$, 由 \eqref{eq:5.63} 还有
$\int_{\mathcal{C}_{-\ell,\ell}}|v_\theta^j|^2>\delta>0$. 此外, $v^j$
也是 $1/j$-几乎调和的(这一性质在伸缩下不变), 故仍可像
引理~\ref{lem:5.16} 中那样论证, 得到子列在
$\mathcal{C}_{-3\ell,3\ell}$ 的紧子集上 $W^{1,2}$ 收敛到调和映射
$v:\bS^2\to\R^N$(这里用了 $M_j$ 的某个子列收敛到一个仿射空间).
如前, \eqref{eq:5.63} 蕴含
\begin{equation}
\int_{\mathcal{C}_{-\ell,\ell}}|v_\theta|^2
\geq\delta\int_{\mathcal{C}_{-3\ell,3\ell}}|\nabla v|^2.
\label{eq:5.67}
\end{equation}
这一次归一化给出 $\int_{\mathcal{C}_{-\ell,\ell}}|v_\theta|^2\geq\delta$,
于是 $v$ 与命题~\ref{prop:5.13} 矛盾, 证明完成.
\end{proof}

\mtparasec{命题 \ref{prop:5.18} 的证明}

\begin{proof}[命题~\ref{prop:5.18}的证明]
对每个整数 $j=0,\dots,m$, 令
$\mathcal{C}(j)=\mathcal{C}_{-(j+3)\ell,(3-j)\ell}$, 并取
引理~\ref{lem:5.19} 给出的 $\mu>0$. 若 $u$ 在 $\mathcal{C}(j)$ 上的限制
是 $\mu$-几乎调和的, 则称第 $j$ 个柱面 $\mathcal{C}(j)$ 是\textbf{好的};
否则称 $\mathcal{C}(j)$ 是\textbf{坏的}.

在每个好的 $\mathcal{C}(j)$ 上应用引理~\ref{lem:5.19}, 得
\begin{equation}
\int_{\mathcal{C}_{-(j+1)\ell,(1-j)\ell}}|u_\theta|^2
\leq\delta\int_{\mathcal{C}(j)}|\nabla u|^2,
\label{eq:5.68}
\end{equation}
于是对好的 $j$ 求和给出
\begin{equation}
\sum_{j\ \text{好}}\int_{\mathcal{C}_{-(j+1)\ell,(1-j)\ell}}|u_\theta|^2
\leq\delta\sum_{j\ \text{好}}\int_{\mathcal{C}(j)}|\nabla u|^2
\leq 6\delta\int_{\mathcal{C}_{-(m+3)\ell,3\ell}}|\nabla u|^2,
\label{eq:5.69}
\end{equation}
其中最后一个不等式用了每个 $\mathcal{C}_{i,i+1}$ 至多含于 $6$ 个
$\mathcal{C}(j)$ 之中.

我们将通过证明坏的 $\mathcal{C}(j)$ 上的总能量(而不只是 $\theta$-能量)
很小来完成证明. 按定义, 对每个坏的 $\mathcal{C}(j)$, 可以选取
$\mathcal{C}(j)$ 中一族有限个互不相交的闭球 $\mathcal{B}_j$, 使得若
$v:\frac18\mathcal{B}_j\to M$ 是任一在
$\partial\frac18\mathcal{B}_j$ 上等于 $u$ 的极小化能量映射, 则
\begin{equation}
\int_{\frac18\mathcal{B}_j}|\nabla u-\nabla v|^2\geq a_j
>\mu\int_{\mathcal{C}(j)}|\nabla u|^2.
\label{eq:5.70}
\end{equation}
由于每个 $\mathcal{C}(j)$ 的内部只与满足 $0<|j-k|\leq 5$ 的
$\mathcal{C}(k)$ 相交, 我们可以把坏的 $\mathcal{C}(j)$ 分成十个子族, 使得
每个子族中各 $\mathcal{C}(j)$ 的内部互不相交. 特别地, 其中某个互不相交的
子族, 记作 $\Gamma$, 满足
\begin{equation}
\sum_{j\in\Gamma}a_j\geq\frac1{10}\sum_{j\ \text{坏}}a_j
\geq\frac1{10}\sum_{j\ \text{坏}}\mu\int_{\mathcal{C}(j)}|\nabla u|^2,
\label{eq:5.71}
\end{equation}
其中最后一个不等式用了 \eqref{eq:5.70}.

然而, 由于 $\bigcup_{j\in\Gamma}\mathcal{B}_j$ 本身就是整个柱面
$\mathcal{C}_{-(m+3)\ell,3\ell}$ 中一族有限个互不相交的闭球, 而 $u$ 在
$\mathcal{C}_{-(m+3)\ell,3\ell}$ 上是 $\nu$-几乎调和的, 故有
\begin{equation}
\frac{\mu}{10}\sum_{j\ \text{坏}}\int_{\mathcal{C}(j)}|\nabla u|^2
\leq\nu\int_{\mathcal{C}_{-(m+3)\ell,3\ell}}|\nabla u|^2.
\label{eq:5.72}
\end{equation}
为得到命题, 把 \eqref{eq:5.69} 与 \eqref{eq:5.72} 结合起来, 得
\begin{equation}
\int_{\mathcal{C}_{-m\ell,0}}|u_\theta|^2
\leq\Bigl(6\delta+\frac{10\nu}{\mu}\Bigr)
\int_{\mathcal{C}_{-(m+3)\ell,3\ell}}|\nabla u|^2.
\label{eq:5.73}
\end{equation}
最后, 取 $\nu$ 充分小即完成证明.
\end{proof}

\mtpara{气泡紧性}
下面用 \cite{PkW} 中为伪全纯曲线发展、后被 \cite{Pk} 用于调和映射的
重整化程序的变体来证明命题~\ref{prop:5.12}. 证明的关键在于能量的一致界、
(A) 与 (B) 都在伸缩下不变, 因而对 $u^j$ 与 $\bS^2$ 的任意一列共形伸缩
的复合也适用.

\begin{proof}[命题~\ref{prop:5.12}的证明]
我们将利用能量界与 (B) 证明 $u^j$ 的某个子列按\defref{def:3.26} 的 (B1)、(B2) 与
(B3) 的意义收敛到一族调和映射. 然后回过头来利用 (A) 与 (B) 证明能量等式
(B4) 也成立. 于是该子列气泡收敛, 从而由\propref{prop:3.27} 也变分流形收敛.

取 $\delta=1/21$, 并取命题~\ref{prop:5.13} 给出的 $\ell\geq 1$ 与
$\epsilon_2>0$. 令 $\epsilon_3=\min\{\epsilon_0/2,\epsilon_2\}$.

\textbf{第一步: 初始紧性.} 引理~\ref{lem:5.16} 给出有限奇异点集
$\mathcal{S}_0\subset\bS^2$、调和映射 $v_0:\bS^2\to M$ 以及子列(仍记为
$u^j$), 它在 $W^{1,2}(\bS^2)$ 中弱收敛到 $v_0$, 且对任一紧子集
$K\subset\bS^2\setminus\mathcal{S}_0$ 在 $W^{1,2}(K)$ 中强收敛. 此外,
测度 $|\nabla u^j|^2dx$ 收敛到测度 $\nu_0$, 满足
$\nu_0(\bS^2)\leq E_0$, 且每个奇异点 $x\in\mathcal{S}_0$ 都有
$\nu_0(x)\geq\epsilon_0$.

\textbf{第二步: 在奇异点处重整化.} 设 $x\in\mathcal{S}_0$ 是第一步得到的
奇异点. 取半径 $\rho>0$, 使得 $x$ 是 $B_{2\rho}(x)$ 中唯一的奇异点, 且
$\int_{B_\rho(x)}|\nabla v_0|^2<\epsilon_3/3$. 对每个 $j$, 令 $r_j>0$ 为
使
\begin{equation}
\inf_{y\in B_{\rho-r_j}(x)}\int_{B_\rho(x)\setminus B_{r_j}(y)}
|\nabla u^j|^2=\epsilon_3
\label{eq:5.74}
\end{equation}
成立的最小半径, 并取球 $B_{r_j}(y_j)\subset B_\rho(x)$, 使得
\begin{equation}
\int_{B_\rho(x)\setminus B_{r_j}(y_j)}|\nabla u^j|^2=\epsilon_3.
\label{eq:5.75}
\end{equation}
由于 $u^j$ 在 $B_\rho(x)\setminus\{x\}$ 的紧子集上收敛到 $v_0$, 故
$y_j\to x$ 且 $r_j\to 0$. 对每个 $j$, 令
\begin{equation}
\Psi_j:\R^2\to\R^2
\label{eq:5.76}
\end{equation}
为把 $B_{r_j}(y_j)$ 送到单位球 $B_1(0)\subset\R^2$ 的``伸缩''. 由伸缩下
的不变性, 伸缩后的映射
$\tilde{u}_1^j=u^j\circ\Psi_j^{-1}$ 仍满足 (B) 且能量不变. 于是
引理~\ref{lem:5.16} 给出子列(仍记为 $\tilde{u}_1^j$)、有限奇异集
$\mathcal{S}_1$ 以及调和映射 $v_1$, 使得 $\tilde{u}_1^j\circ\Pi$ 在
$W^{1,2}(\bS^2)$ 中弱收敛到 $v_1$, 且对任一紧子集
$K\subset\bS^2\setminus\mathcal{S}_1$ 在 $W^{1,2}(K)$ 中强收敛. 此外,
测度 $|\nabla\tilde{u}_1^j\circ\Pi|^2dx$ 收敛到测度 $\nu_1$.

球 $B_{r_j}(y_j)$ 的选取保证
$\nu_1(\bS^2\setminus\{p^+\})\leq\nu_0(x)$ 以及
$\nu_1(S^-)\leq\nu_0(x)-\epsilon_3$. (回忆球极投影 $\Pi$ 把开南半球
$S^-$ 送到 $\R^2$ 中的开单位球.) 迭代的关键是下述断言:

\begin{quote}
$(\star)$ 在任一 $y\in\mathcal{S}_1\setminus\{p^+\}$ 处的最大能量集中
不超过 $\nu_0(x)-\epsilon_3/3$.
\end{quote}

由于奇异点处的能量或非平凡调和映射的能量至少为
$\epsilon_0>\epsilon_3$, $(\star)$ 可能不成立的唯一方式是: $v_1$ 为常值,
$\mathcal{S}_1$ 恰为两点 $p^+$ 与 $y$, 且 $\nu_0(x)$ 中至多
$\epsilon_3/3$ 的部分在 $p^+$ 处逸出. 但这样将蕴含
$\int_{B_\rho(x)}|\nabla u^j|^2$ 中除至多 $2\epsilon_3/3$ 之外全部位于
$B_{t_j}(y_j)$ 中, 而 $t_j/r_j\to 0$, 这与 $r_j$ 的极小性矛盾.

\textbf{第三步: 重复上述过程.} 我们对 $\mathcal{S}_0$ 中其余的奇异点、
以及南半球中的每个奇异点 $\mathcal{S}_1$ 等等重复这一爆破构造, 以得到新的
极限调和映射与新的待爆破奇异点. 由 $(\star)$ 知此过程必在至多
$3E_0/\epsilon_3$ 步后终止.

\textbf{第四步: 颈部.} 我们已证明 $u^j$ 按 (B1)、(B2) 与 (B3) 的意义
收敛到一族调和映射. 还需证明 (B4), 即 $v_k$ 们解释了序列 $u^j$ 的全部
能量, 极限中没有能量损失.

为理解能量可能如何损失, 重新考察爆破过程中能量的去向是有益的. 在爆破
过程的每一阶段, 能量从某个奇异点 $x$ ``取出'', 然后去往两个地方之一:

\begin{itemize}[leftmargin=1.8em,itemsep=1pt,topsep=2pt]
\item 它可能出现在 $\bS^2\setminus\{p^+\}$ 中某个奇异点的新极限调和
      映射里.
\item 它可能在北极 $p^+$ 处消失(即
      $\nu_1(\bS^2\setminus\{p^+\})<\nu_0(x)$).
\end{itemize}

第一种情形下, 能量在极限中被解释清楚, 或留存到更晚的阶段. 然而第二种
情形下, 能量永久丢失, 这正是我们必须排除的.

我们用反证法, 设
\begin{equation}
\nu_1\bigl(\bS^2\setminus\{p^+\}\bigr)<\nu_0(x)-\hat{\delta}
\label{eq:5.77}
\end{equation}
对某个 $\hat{\delta}>0$ 成立. (注意必有 $\hat{\delta}\leq\epsilon_3$.)
沿用第一步的记号, 于是设 $A_j=B_{s_j}(y_j)\setminus B_{t_j}(y_j)$ 是
圆环, 满足
\begin{equation}
s_j\to 0,\qquad \frac{t_j}{r_j}\to\infty,\qquad
\int_{A_j}|\nabla u_j|^2\geq\hat{\delta}>0.
\label{eq:5.78}
\end{equation}
显然 $s_j$ 与 $t_j$ 的选取有相当大的自由度. 特别地, 可以选取一列
$\lambda_j\to\infty$, 使得圆环
\begin{equation}
\tilde{A}_j=B_{\rho/2}(y_j)\setminus B_{t_j/\lambda_j}(y_j)
\label{eq:5.79}
\end{equation}
也满足这一点, 即 $\lambda_js_j\to 0$ 且
$t_j/(\lambda_jr_j)\to\infty$. 由 \eqref{eq:5.78} 与 $r_j$ 的定义推出
\begin{equation}
\int_{\tilde{A}_j}|\nabla u^j|^2\leq\epsilon_3\leq\epsilon_2.
\label{eq:5.80}
\end{equation}
然而, 把它与命题~\ref{prop:5.18}(取 $\delta=1/21$)结合起来便表明: 当
$j$ 大时面积必严格小于能量, 与 (A) 矛盾, 从而完成证明.
\end{proof}

% ================================================================
\section{极小球面与宽度}
\label{sec:5.5}

下面定义一个二维的宽度, 它与前面定义的一维宽度类似, 并用它来产生极小球面.

\mtpara{宽度}
现在给出前面为曲线扫掠定义的宽度的一个二维版本. 我们有时称之为
$2$-宽度, 但在不会引起混淆时, 就简单地称之为宽度. 记 $\Omega$ 为连续
映射
\begin{equation*}
\sigma:\bS^2\times[0,1]\to M
\end{equation*}
之集, 使得:

\begin{itemize}[leftmargin=1.8em,itemsep=1pt,topsep=2pt]
\item 对每个 $t\in[0,1]$, 映射 $\sigma(\cdot,t)$ 属于 $C^0\cap W^{1,2}$.
\item 映射 $t\mapsto\sigma(\cdot,t)$ 是从 $[0,1]$ 到 $C^0\cap W^{1,2}$
      的连续映射.
\item $\sigma$ 把 $\bS^2\times\{0\}$ 与 $\bS^2\times\{1\}$ 映为点.
\end{itemize}

给定映射 $\beta\in\Omega$, 同伦类 $\Omega_\beta$ 定义为 $\Omega$ 中所有与
$\beta$ 在 $\Omega$ 内同伦的映射之集. 我们把任何这样的 $\beta$ 称为一个
\newterm{sweepout}{扫掠}{sweepout}; 有些作者采用更严格的定义, 要求
$\beta$ 还必须诱导出从 $\bS^3$ 到 $M$ 的映射度为 $1$ 的映射. 事实上,
我们最感兴趣的是 $\beta$ 诱导出从 $\bS^3$ 到 $M$ 的、属于 $\pi_3(M)$ 中
非平凡类\footnote{例如, 当 $M$ 是同伦 $3$-球面且诱导映射的映射度为 $1$
时.}的映射的情形. 理由是在这种情形下宽度为正, 并且正如我们将看到的,
它等于一族非空的极小 $2$-球面的面积.

与同伦类 $\Omega_\beta$ 相联系的(\textbf{能量})宽度
$W_E=W_E(\beta,M)$ 定义为各切片能量极大值之下确界. 也就是说, 令
\begin{equation}
W_E=\inf_{\sigma\in\Omega_\beta}\max_{t\in[0,1]}
\mathrm{E}\bigl(\sigma(\cdot,t)\bigr),
\label{eq:5.81}
\end{equation}
其中能量由
\begin{equation}
\mathrm{E}\bigl(\sigma(\cdot,t)\bigr)
=\frac12\int_{\bS^2}|\nabla_x\sigma(x,t)|^2dx
\label{eq:5.82}
\end{equation}
给出.

尽管这类构造总被称为极小极大, 它实际上是下确界--极大. 也就是说, 对每个
(光滑)扫掠先看各切片能量的极大值, 再对给定同伦类中所有扫掠取下确界.
按定义宽度总是非负的, 而当 $\beta$ 的同伦类非平凡时为正. 例如, 正性可以
直接用 \cite{Jo} 看出. 确切地说, \cite{Jo} 第 125 页证明: 若
$\max_t\mathrm{E}(\sigma(\cdot,t))$ 充分小(依赖于 $M$), 则 $\sigma$ 是
同伦平凡的.\footnote{推论~\ref{cor:5.10} 可用来给出另一个证明.}

人们也可以改用面积而非能量来定义宽度, 即令
\begin{equation}
W_A=\inf_{\sigma\in\Omega_\beta}\max_{t\in[0,1]}
\Area\bigl(\sigma(\cdot,t)\bigr).
\label{eq:5.83}
\end{equation}
$W^{1,2}$ 映射 $u:\bS^2\to\R^N$ 的面积按定义是 Jacobi 行列式的积分
\begin{equation}
J_u=\sqrt{\det\bigl(du^Tdu\bigr)},
\label{eq:5.84}
\end{equation}
其中 $du$ 是 $u$ 的微分, $du^T$ 是它的转置. 也就是说, 若 $e_1,e_2$ 是
$D\subset\bS^2$ 上的正交标架, 则
\begin{equation}
J_u=\Bigl(|u_{e_1}|^2|u_{e_2}|^2-\langle u_{e_1},u_{e_2}\rangle^2
\Bigr)^{\frac12}\leq\frac12|du|^2
\label{eq:5.85}
\end{equation}
并且
\begin{equation}
\Area\bigl(u\big|_D\bigr)=\int_DJ_u\leq\mathrm{E}\bigl(u\big|_D\bigr).
\label{eq:5.86}
\end{equation}
因此, 面积小于或等于能量, 且等号成立当且仅当
$\langle u_{e_1},u_{e_2}\rangle$ 与 $|u_{e_1}|^2-|u_{e_2}|^2$ 作为
$L^1$ 函数为零. 等号成立时, 称 $u$ 是\termcn{几乎共形}%
（almost conformal）的; 见 \eqref{eq:5.86}. 与经典 Plateau 问题一样
(见\thmref{thm:4.1}), 在证明极小曲面的存在性时能量用起来更方便些.
下一个命题表明, 与 Plateau 问题一样有 $W_E=W_A$(由上面的讨论显然
$W_A\leq W_E$). 因此我们将略去下标, 直接写作 $W$.

\begin{proposition}
\label{prop:5.20}
(Colding--Minicozzi, \cite{CM27}) $W_E=W_A$.
\end{proposition}

由 \eqref{eq:5.86}, 命题~\ref{prop:5.20} 在证明了 $W_E\leq W_A$ 之后
便得. Plateau 问题的相应结果是通过取面积的极小化序列并重新参数化使这些
映射共形(即选取等温坐标)来证明的. 具体实施时有几点技术困难, 因为拉回
度量可能是退化的, 而且只在 $L^1$ 中, 而等温坐标的存在性要求诱导度量为正
且有界; 例如见 \cite{ScW} 的命题 5.4.

我们这里沿用同样的途径, 区别在于需要让重新参数化连续依赖于 $t$.

\mtpara{光滑映射的稠密性}
下一个引理注意到, Schoen 与 Uhlenbeck 在 \cite{ScU2} 第 4 节命题中所用的
卷积正则化是连续的.

\begin{lemma}
\label{lem:5.21}
给定 $\gamma\in\Omega$ 与 $\epsilon>0$, 存在正则化 $\tilde{\gamma}\in
\Omega_\gamma$, 使得
\begin{equation}
\max_t\|\tilde{\gamma}(\cdot,t)-\gamma(\cdot,t)\|_{W^{1,2}}\leq\epsilon,
\label{eq:5.87}
\end{equation}
每个切片 $\tilde{\gamma}(\cdot,t)$ 都是 $C^2$ 的, 且映射
$t\mapsto\tilde{\gamma}(\cdot,t)$ 是从 $[0,1]$ 到
$C^2(\bS^2,M)$ 的连续映射.
\end{lemma}

\begin{proof}
由于 $M$ 光滑、紧且嵌入, 存在 $\delta>0$, 使得对 $M$ 在 $\R^N$ 中的
$\delta$-管状邻域 $M_\delta$ 中的每个 $x$, 都有唯一的最近点
$\Pi(x)\in M$, 且映射 $x\mapsto\Pi(x)$ 光滑. $\Pi$ 称为\newterm{nearest}{最近点投影}{nearest point projection}.

给定开球 $B_1(0)\subset\R^3$ 中的 $y$, 定义 $T_y:\bS^2\to\bS^2$ 为
$T_y(x)=\frac{x-y}{|x-y|}$. 由于每个 $T_y$ 光滑且这些映射光滑依赖于
$y$, 映射 $(y,f)\mapsto f\circ T_y$ 是连续的:
\begin{equation}
B_1(0)\times C^0\cap W^{1,2}(\bS^2,\R^N)
\to C^0\cap W^{1,2}(\bS^2,\R^N)
\label{eq:5.88}
\end{equation}
(这对 $f\in C^1$ 是显然的, 对 $C^0\cap W^{1,2}$ 则由稠密性推出). 因此,
由于 $T_0$ 是恒等映射, 给定 $f\in C^0\cap W^{1,2}(\bS^2,\R^N)$ 与
$\mu>0$, 存在 $r>0$, 使得
\begin{equation}
\sup_{|y|\leq r}\|f\circ T_y-f\|_{C^0\cap W^{1,2}}<\mu.
\label{eq:5.89}
\end{equation}
把上式用于每个 $t$ 对应的 $\gamma(\cdot,t)$, 并利用
$t\mapsto\gamma(\cdot,t)$ 到 $C^0\cap W^{1,2}$ 连续以及 $[0,1]$ 紧,
得到 $\bar{r}>0$, 使得
\begin{equation}
\sup_{t\in[0,1]}\sup_{|y|\leq\bar{r}}
\|T_y\gamma(\cdot,t)-\gamma(\cdot,t)\|_{C^0\cap W^{1,2}}<\mu.
\label{eq:5.90}
\end{equation}

其次, 固定一个光滑的径向磨光核 $\phi\geq 0$, 其积分为 $1$, 支集含于
$\R^3$ 的单位球中. 对每个 $r\in(0,1)$, 定义
$\phi_r(x)=r^{-3}\phi(x/r)$, 并令
\begin{equation}
\gamma_r(x,t)=\int_{B_r(0)}\phi_r(y)\gamma\bigl(T_y(x),t\bigr)dy
=\int_{B_r(x)}\phi_r(x-y)\gamma\Bigl(\frac{y}{|y|},t\Bigr)dy.
\label{eq:5.91}
\end{equation}
关于磨光核卷积, 我们有下述标准性质(例如见 \cite{Ev} 的 5.3 节与附录
C.4): 首先, 每个 $\gamma_r(\cdot,t)$ 光滑, 且对每个 $k$, 映射
$t\mapsto\gamma_r(\cdot,t)$ 是从 $[0,1]$ 到 $C^k(\bS^2,\R^N)$ 的连续
映射. 其次,
\begin{equation}
\begin{aligned}
\|\gamma_r(\cdot,t)-\gamma(\cdot,t)\|_{C^0}^2
&\leq\sup_{|y|\leq r}\|T_y\gamma(\cdot,t)-\gamma(\cdot,t)\|_{C^0}^2,\\
\|\nabla\gamma_r(\cdot,t)-\nabla\gamma(\cdot,t)\|_{L^2}^2
&\leq\sup_{|y|\leq r}\|T_y\gamma(\cdot,t)-\gamma(\cdot,t)\|_{L^2}^2.
\end{aligned}
\label{eq:5.92}
\end{equation}
由 \eqref{eq:5.92} 与 \eqref{eq:5.90} 推出, 当 $r\leq\bar{r}$ 时对一切
$t$ 有
\begin{equation}
\|\gamma_r(\cdot,t)-\gamma(\cdot,t)\|_{C^0\cap W^{1,2}}<\mu.
\label{eq:5.93}
\end{equation}
映射 $\gamma_r(\cdot,t)$ 未必落在 $M$ 中, 但当 $\mu$ 小时由
\eqref{eq:5.93} 它落在 $M_\delta$ 中. 于是映射
$\tilde{\gamma}(x,t)=\Pi\circ\gamma_r(x,t)$ 满足 \eqref{eq:5.87}, 每个
切片 $\tilde{\gamma}(\cdot,t)$ 是 $C^2$ 的, 且
$t\mapsto\tilde{\gamma}(\cdot,t)$ 是从 $[0,1]$ 到
$C^2(\bS^2,M)$ 的连续映射. 最后, $s\mapsto\tilde{\gamma}_{sr}$ 是连接
$\tilde{\gamma}$ 与 $\gamma$ 的一个显式同伦.
\end{proof}

\mtpara{能量与面积的等价性}
我们还需要等温坐标的存在性, 并特别留意其对度量的依赖. 记
$\bS^2_{g_0}$ 为 $\bS^2$ 上常曲率 $1$ 的圆度量.

\begin{lemma}
\label{lem:5.22}
给定 $\bS^2$ 上的 $C^1$ 度量 $\tilde{g}$, 存在唯一的保向
$C^{1,1/2}$ 共形微分同胚
$h_{\tilde{g}}:\bS^2_{g_0}\to\bS^2_{\tilde{g}}$, 它固定给定的 $3$ 个点.

此外, 若 $\tilde{g}_1$ 与 $\tilde{g}_2$ 是两个 $C^1$ 度量, 且都存在
$\epsilon>0$ 使二者都 $\geq\epsilon g_0$, 则
\begin{equation}
\|h_{\tilde{g}_1}-h_{\tilde{g}_2}\|_{C^0\cap W^{1,2}}
\leq C\|\tilde{g}_1-\tilde{g}_2\|_{C^0},
\label{eq:5.94}
\end{equation}
其中常数 $C$ 依赖于 $\epsilon$ 与 $\tilde{g}_i$ 的 $C^1$ 范数的最大值.
\end{lemma}

\begin{proof}
变度量情形的 Riemann 映射定理(见 \cite{Jo} 的定理 3.1.1 与推论 3.1.1;
参看 \cite{ABe} 或 \cite{Mo2})给出共形微分同胚
$h_{\tilde{g}}:\bS^2_{g_0}\to\bS^2_{\tilde{g}}$.

我们将分别界定 $C^0$ 范数与 $W^{1,2}$ 范数. 首先, \cite{ABe} 的引理 17
给出
\begin{equation}
\|h_{\tilde{g}_1}-h_{\tilde{g}_2}\|_{C^0}
\leq C_1\|\tilde{g}_1-\tilde{g}_2\|_{C^0},
\label{eq:5.95}
\end{equation}
其中 $C_1$ 依赖于 $\epsilon$ 与这些度量的 $C^0$ 范数. 其次,
\cite{ABe} 的定理 8 对 $\bS^2$ 中任一单位球上
$\nabla(h_{\tilde{g}_1}-h_{\tilde{g}_2})$ 给出一个一致的 $L^p$ 界,
其中由 \cite{ABe} 的 (8) 有 $p>2$:
\begin{equation}
\|\nabla(h_{\tilde{g}_1}-h_{\tilde{g}_2})\|_{L^p(B_1)}
\leq C_2\|\tilde{g}_1-\tilde{g}_2\|_{C^0(\bS^2)},
\label{eq:5.96}
\end{equation}
其中 $C_2$ 依赖于 $\epsilon$ 与这些度量的 $C^0$ 范数. 用有限个这样的单位
球覆盖 $\bS^2$ 并应用 H\"older 不等式, 即得所需的能量界.
\end{proof}

现在可以证明两个宽度相等.

\begin{proof}[命题~\ref{prop:5.20}的证明]
由 \eqref{eq:5.86} 有 $W_A\leq W_E$. 为证明 $W_E\leq W_A$, 给定
$\epsilon>0$, 取扫掠 $\gamma\in\Omega_\beta$ 使
\begin{equation}
\max_{t\in[0,1]}\Area\bigl(\gamma(\cdot,t)\bigr)<W_A+\epsilon/2.
\label{eq:5.97}
\end{equation}
由引理~\ref{lem:5.21}, 存在正则化 $\tilde{\gamma}\in\Omega_\beta$,
使每个切片 $\tilde{\gamma}(\cdot,t)$ 为 $C^2$, 映射
$t\mapsto\tilde{\gamma}(\cdot,t)$ 从 $[0,1]$ 到 $C^2(\bS^2,M)$ 连续,
并且(同样由 \eqref{eq:3.112})
\begin{equation}
\max_t\Area\bigl(\tilde{\gamma}(\cdot,t)\bigr)<W_A+\epsilon.
\label{eq:5.98}
\end{equation}
映射 $\tilde{\gamma}(\cdot,t)$ 诱导出 $\bS^2$ 上一族连续的单参数拉回
$C^1$ 度量 $g(t)$(可能退化). 引理~\ref{lem:5.22} 要求度量非退化, 故
定义扰动后的度量 $\tilde{g}(t)=g(t)+\delta g_0$. 对每个 $t$,
引理~\ref{lem:5.22} 给出 $C^{1,1/2}$ 共形微分同胚
$h_t:\bS^2_{g_0}\to\bS^2_{\tilde{g}(t)}$, 它们在
$C^0\cap W^{1,2}(\bS^2,\bS^2)$ 中连续变化. 由此以及
$t\mapsto\tilde{\gamma}(\cdot,t)$ 在 $C^2$ 中的连续性和链式法则, 推出
$t\mapsto\tilde{\gamma}(\cdot,t)\circ h_t$ 作为从 $[0,1]$ 到
$C^0\cap W^{1,2}(\bS^2,M)$ 的映射连续.

下面用 $h_t$ 的共形性, 如 \cite{ScW} 的命题 5.4 那样来控制复合映射的
能量. 也就是说, 我们有
\begin{equation}
\begin{aligned}
\mathrm{E}\bigl(\tilde{\gamma}(\cdot,t)\circ h_t\bigr)
&=\mathrm{E}\bigl(h_t:\bS^2_{g_0}\to\bS^2_{g(t)}\bigr)
\leq\mathrm{E}\bigl(h_t:\bS^2_{g_0}\to\bS^2_{\tilde{g}(t)}\bigr)\\
&=\Area\bigl(\bS^2_{\tilde{g}(t)}\bigr)
=\int_{\bS^2}\bigl[\det\bigl(g_0^{-1}g(t)\bigr)
+\delta\operatorname{Tr}\bigl(g_0^{-1}g(t)\bigr)+\delta^2\bigr]^{1/2}
\,dvol_{g_0}\\
&\leq\Area\bigl(\bS^2_{g(t)}\bigr)
+4\pi\Bigl[\delta^2+2\delta\sup_t\bigl|g_0^{-1}g(t)\bigr|\Bigr]^{1/2}.
\end{aligned}
\label{eq:5.99}
\end{equation}
取 $\delta>0$ 使
$4\pi[\delta^2+2\delta\sup_t|g_0^{-1}g(t)|]^{1/2}<\epsilon$.

若 $\tilde{\gamma}(\cdot,t)\circ h_t$ 与 $\tilde{\gamma}$ 同伦, 我们便
完成了. 然而, 由 Smale 定理, $\bS^2$ 的保向微分同胚空间同伦等价于
$\mathbb{RP}^3$. 为绕开这一点, 注意
$t\mapsto\|\tilde{\gamma}(\cdot,t)\|_{C^2}$ 连续且在 $t=1$ 处为零, 故
存在 $\tau<1$, 使得
\begin{equation}
\sup_{t\geq\tau}\|\tilde{\gamma}(\cdot,t)\|_{C^2}
\leq\frac{\epsilon}{\sup_{t\in[0,1]}\|h_t\|_{W^{1,2}}^2}.
\label{eq:5.100}
\end{equation}
于是, 若令 $\tilde{h}_t$ 在 $[0,\tau]$ 上等于 $h_t\equiv h(t)$, 而在
$[\tau,1]$ 上等于 $h(\tau(1-t)/(1-\tau))$, 则 \eqref{eq:5.99} 与
\eqref{eq:5.100} 蕴含
\begin{equation}
\max_{t\in[0,1]}\mathrm{E}\bigl(\tilde{\gamma}(\cdot,t)\circ\tilde{h}_t\bigr)
\leq W_A+2\epsilon.
\label{eq:5.101}
\end{equation}
此外, 映射 $\tilde{\gamma}(\cdot,t)\circ\tilde{h}_t$ 也属于 $\Omega$.
最后, 把 $\tau$ 换为 $s\tau$ 并令 $s\to0$, 便给出 $\Omega$ 中从
$\tilde{\gamma}(\cdot,t)\circ\tilde{h}_t$ 到 $\tilde{\gamma}(\cdot,t)$
的一个显式同伦.
\end{proof}

\mtpara{扫掠上的能量下降映射}
现在调和替换的基本性质已经建立, 我们来定义一个从 $\Omega$ 到自身的能量
下降映射, 它保持同伦类(即把每个 $\Omega_\beta$ 映入自身), 并记录其关键
性质. 这应看作 Birkhoff 曲线缩短过程的推广, 后者在收紧由曲线构成的扫掠
时起类似的作用.

\begin{theorem}
\label{thm:5.23}
(Colding--Minicozzi, \cite{CM27}) 存在常数 $\epsilon_0>0$ 与连续函数
$\Psi:[0,\infty)\to[0,\infty)$, $\Psi(0)=0$, 二者都依赖于 $M$, 使得
给定任一无非常值调和切片的 $\tilde{\gamma}\in\Omega$ 以及 $W>0$, 存在
$\gamma\in\Omega_{\tilde{\gamma}}$, 使得对每个 $t$ 有
$\mathrm{E}(\gamma(\cdot,t))\leq\mathrm{E}(\tilde{\gamma}(\cdot,t))$,
并且对每个满足 $\mathrm{E}(\tilde{\gamma}(\cdot,t))\geq W/2$ 的 $t$:

\medskip
\noindent$(B_\Psi)$ 若 $\mathcal{B}$ 是 $\bS^2$ 中任一由有限个互不
相交的闭球构成的族, 且
$\int_{\bigcup_{\mathcal{B}}B}|\nabla\gamma(\cdot,t)|^2<\epsilon_0$, 而
$v:\bigcup_{\mathcal{B}}\frac18B\to M$ 是在
$\bigcup_{\mathcal{B}}\frac18\partial B$ 上等于 $\gamma(\cdot,t)$ 的
极小化能量映射, 则
\begin{equation*}
\int_{\bigcup_{\mathcal{B}}\frac18B}
\bigl|\nabla\gamma(\cdot,t)-\nabla v\bigr|^2
\leq\Psi\bigl[\mathrm{E}(\tilde{\gamma}(\cdot,t))
-\mathrm{E}(\gamma(\cdot,t))\bigr].
\end{equation*}
\end{theorem}

我们将通过对 $\bS^2$ 中一族随 $t$ 连续变化的球作调和替换, 由
$\tilde{\gamma}(\cdot,t)$ 构造出 $\gamma(\cdot,t)$. 这些球将在下文
引理~\ref{lem:5.25} 中选取. 本小节中, $\epsilon_1>0$ 始终表示
定理~\ref{thm:5.8} 给出的(依赖于 $M$ 的)小能量常数.

给定 $\sigma\in\Omega$ 与 $\epsilon\in(0,\epsilon_1]$, 定义在能量至多
$\epsilon$ 的球族上作调和替换所能带来的最大改进为
\begin{equation}
e_{\sigma,\epsilon}(t)=\sup_{\mathcal{B}}\Bigl\{
\mathrm{E}\bigl(\sigma(\cdot,t)\bigr)
-\mathrm{E}\Bigl(H\bigl(\sigma(\cdot,t),\tfrac12\mathcal{B}\bigr)\Bigr)
\Bigr\},
\label{eq:5.102}
\end{equation}
其中上确界取遍所有由互不相交的闭球构成、且 $\sigma(\cdot,t)$ 在其上总
能量至多 $\epsilon$ 的有限族 $\mathcal{B}$. 注意
$e_{\sigma,\epsilon}(t)$ 非负, 关于 $\epsilon$ 单调不减, 且当
$\sigma(\cdot,t)$ 不调和时为正.

\begin{lemma}
\label{lem:5.24}
若 $\sigma(\cdot,t)$ 不调和且 $\epsilon\in(0,\epsilon_1]$, 则存在含
$t$ 的开区间 $I^t$, 使得对二倍区间 $2I^t$ 中的所有 $s$ 有
$e_{\sigma,\epsilon/2}(s)\leq 2e_{\sigma,\epsilon}(t)$.
\end{lemma}

\begin{proof}
由推论~\ref{cor:5.10} 中的 \eqref{eq:5.35}, 存在 $\delta_1>0$(依赖于
$t$), 使得若
\begin{equation}
\|\sigma(\cdot,t)-\sigma(\cdot,s)\|_{C^0\cap W^{1,2}}<\delta_1
\label{eq:5.103}
\end{equation}
且 $\mathcal{B}$ 是一族互不相交的闭球, $\sigma(\cdot,t)$ 与
$\sigma(\cdot,s)$ 在其上的能量都不超过 $\epsilon_1$, 则
\begin{equation}
\Bigl|\mathrm{E}\Bigl(H\bigl(\sigma(\cdot,s),\tfrac12\mathcal{B}\bigr)\Bigr)
-\mathrm{E}\Bigl(H\bigl(\sigma(\cdot,t),\tfrac12\mathcal{B}\bigr)\Bigr)\Bigr|
\leq e_{\sigma,\epsilon}(t)/2.
\label{eq:5.104}
\end{equation}
这里用了 $\sigma(\cdot,t)$ 不调和, 故 $e_{\sigma,\epsilon}(t)>0$. 由于
$t\mapsto\sigma(\cdot,t)$ 作为到 $C^0\cap W^{1,2}$ 的映射连续, 可以
选取 $I^t$, 使得对所有 $s\in 2I^t$ 有 \eqref{eq:5.103} 成立, 并且
\begin{equation}
\frac12\int_{\bS^2}\Bigl||\nabla\sigma(\cdot,t)|^2
-|\nabla\sigma(\cdot,s)|^2\Bigr|
\leq\min\Bigl\{\frac{\epsilon}{2},\frac{e_{\sigma,\epsilon}(t)}{2}\Bigr\}.
\label{eq:5.105}
\end{equation}

现设 $s\in 2I^t$, 且 $\sigma(\cdot,s)$ 在一族球 $\mathcal{B}$ 上的能量
至多 $\epsilon/2$. 由 \eqref{eq:5.105} 知 $\sigma(\cdot,t)$ 在
$\mathcal{B}$ 上的能量至多 $\epsilon$. 结合 \eqref{eq:5.104} 与
\eqref{eq:5.105} 得
\begin{equation}
\begin{aligned}
\Bigl|&\mathrm{E}\bigl(\sigma(\cdot,s)\bigr)
-\mathrm{E}\Bigl(H\bigl(\sigma(\cdot,s),\tfrac12\mathcal{B}\bigr)\Bigr)
-\mathrm{E}\bigl(\sigma(\cdot,t)\bigr)\\
&+\mathrm{E}\Bigl(H\bigl(\sigma(\cdot,t),\tfrac12\mathcal{B}\bigr)\Bigr)
\Bigr|\leq e_{\sigma,\epsilon}(t).
\end{aligned}
\label{eq:5.106}
\end{equation}
由于这对任何这样的 $\mathcal{B}$ 都成立, 我们得到
$e_{\sigma,\epsilon/2}(s)\leq 2e_{\sigma,\epsilon}(t)$.
\end{proof}

对无调和切片的扫掠, 下一个引理构造出有限多个球族, 使得在其中至少一个
球族上作调和替换会严格降低能量. 此外, 每个球族都由有限个两两不相交的
闭球构成.

\begin{lemma}
\label{lem:5.25}
若 $W>0$ 且 $\tilde{\gamma}\in\Omega$ 没有非常值的调和切片, 则得到
整数 $m$(依赖于 $\tilde{\gamma}$)、$m$ 个 $\bS^2$ 中的球族
$\mathcal{B}_1,\dots,\mathcal{B}_m$(每个球族 $\mathcal{B}_j$ 中的球
两两不相交)以及 $m$ 个连续函数
$r_1,\dots,r_m:[0,1]\to[0,1]$, 使得对每个 $t$:

\begin{enumerate}
\item 至多有两个 $r_j(t)$ 为正, 且对每个 $j$ 有
      $\sum_{B\in\mathcal{B}_j}\frac12\int_{r_j(t)B}
      |\nabla\tilde{\gamma}(\cdot,t)|^2<\epsilon_1/3$.
\item {\setlength{\emergencystretch}{2em}%
      若 $\mathrm{E}(\tilde{\gamma}(\cdot,t))\geq W/2$, 则存在 $j(t)$,
      使得在 $\frac{r_{j(t)}}2\mathcal{B}_{j(t)}$ 上作调和替换使能量
      至少下降 $e_{\tilde{\gamma},\epsilon_1/8}(t)/8$.\par}
\end{enumerate}
\end{lemma}

\begin{proof}
由于切片的能量关于 $t$ 连续, 集合
\begin{equation}
I=\{t\mid\mathrm{E}(\tilde{\gamma}(\cdot,t))\geq W/2\}
\label{eq:5.107}
\end{equation}
是紧的. 对每个 $t\in I$, 选取 $\bS^2$ 中一族有限个互不相交的闭球
$\mathcal{B}^t$, 使
$\frac12\int_{\bigcup_{\mathcal{B}^t}}|\nabla\tilde{\gamma}(\cdot,t)|^2
\leq\epsilon_1/4$, 于是
\begin{equation}
\mathrm{E}\bigl(\gamma(\cdot,t)\bigr)
-\mathrm{E}\Bigl(H\bigl(\gamma(\cdot,t),\tfrac12\mathcal{B}^t\bigr)\Bigr)
\geq\frac{e_{\tilde{\gamma},\epsilon_1/4}(t)}{2}>0.
\label{eq:5.108}
\end{equation}
引理~\ref{lem:5.24} 给出含 $t$ 的开区间 $I^t$, 使得对所有 $s\in 2I^t$
有
\begin{equation}
e_{\tilde{\gamma},\epsilon_1/8}(s)\leq 2e_{\tilde{\gamma},\epsilon_1/4}(t).
\label{eq:5.109}
\end{equation}
利用 $\tilde{\gamma}(\cdot,s)$ 在 $C^0\cap W^{1,2}$ 中的连续性与推论
~\ref{cor:5.10}, 可以缩小 $I^t$, 使得当 $s\in 2I^t$ 时
$\tilde{\gamma}(\cdot,s)$ 在 $\mathcal{B}^t$ 上的能量至多 $\epsilon_1/3$,
并且还有
\begin{equation}
\begin{aligned}
\Bigl|&\mathrm{E}\bigl(\gamma(\cdot,s)\bigr)
-\mathrm{E}\Bigl(H\bigl(\gamma(\cdot,s),\tfrac12\mathcal{B}^t\bigr)\Bigr)
-\mathrm{E}\bigl(\gamma(\cdot,t)\bigr)\\
&+\mathrm{E}\Bigl(H\bigl(\gamma(\cdot,t),\tfrac12\mathcal{B}^t\bigr)\Bigr)
\Bigr|\leq\frac{e_{\tilde{\gamma},\epsilon_1/4}(t)}{4}.
\end{aligned}
\label{eq:5.110}
\end{equation}
由于 $I$ 紧, 可用有限个 $I^t$ 覆盖 $I$, 设为
$I^{t_1},\dots,I^{t_m}$. 此外, 在舍弃其中一些区间之后, 可以安排:

\medskip
\noindent\textbf{断言}: 每个 $t$ 至少属于一个闭区间
$\overline{I^{t_j}}$; 每个 $\overline{I^{t_j}}$ 至多与另外两个
$\overline{I^{t_k}}$ 相交; 且与 $\overline{I^{t_j}}$ 相交的那些
$\overline{I^{t_k}}$ 彼此不相交.

\medskip
我们将给出选取这些集合的具体做法来证明这一点. 首先, 若
$\overline{I^{t_1}}$ 含于另外两个区间的并中, 就把它舍去. 否则, 考虑
左端点落在 $\overline{I^{t_1}}$ 中的那些区间, 找出右端点最大的一个, 其余
的(已包含在其中)舍去. 类似地, 考虑右端点落在 $\overline{I^{t_1}}$ 中的
那些区间, 除左端点最小的一个之外全部舍去. 接着对 $I^{t_2}$(若尚未被舍去)
重复此过程, 如此下去. 至多 $m$ 步之后便得到所需的覆盖, 断言得证.

对每个 $j=1,\dots,m$, 选取连续函数 $r_j:[0,1]\to[0,1]$, 使得:

\begin{itemize}[leftmargin=1.8em,itemsep=1pt,topsep=2pt]
\item $r_j(t)=1$ 于 $\overline{I^{t_j}}$ 上, 而当
      $t\notin 2I^{t_j}$ 时 $r_j(t)$ 为零.
\item $r_j(t)$ 在与 $\overline{I^{t_j}}$ 不相交的那些区间上为零.
\end{itemize}

性质 (1) 直接可得, 而 (2) 由 \eqref{eq:5.108}、\eqref{eq:5.109} 与
\eqref{eq:5.110} 推出.
\end{proof}

\begin{proof}[定理~\ref{thm:5.23}的证明]
设 $\mathcal{B}_1,\dots,\mathcal{B}_m$ 与
$r_1,\dots,r_m:[0,1]\to[0,\pi)$ 由引理~\ref{lem:5.25} 给出. 我们将用
一个 $m$ 步替换过程来定义 $\gamma$. 确切地说, 先令
$\gamma^0=\tilde{\gamma}$, 然后对每个 $k=1,\dots,m$, 在第 $k$ 个球族
$r_k(t)\mathcal{B}_k$ 上对 $\gamma^{k-1}(\cdot,t)$ 作调和替换来定义
$\gamma^k$; 也就是说, 令
\begin{equation}
\gamma^k(\cdot,t)=H\bigl(\gamma^{k-1}(\cdot,t),r_k(t)\mathcal{B}_k\bigr).
\label{eq:5.111}
\end{equation}
最后令 $\gamma=\gamma^m$.

构造中的一个要点是: 球族的性质 (1) 表明对每个 $t$ 只有两个 $r_k(t)$ 为正.
因此, 性质 (1) 给出的球上的能量界蕴含每个极小化能量映射所替换的映射能量
至多 $2\epsilon_1/3<\epsilon_1$. 于是推论~\ref{cor:5.10} 蕴含它们连续
依赖于边界值, 而边界值本身关于 $t$ 连续, 故所得映射 $\tilde{\gamma}$ 也
关于 $t$ 连续. 最后, 显然 $\tilde{\gamma}$ 与 $\gamma$ 同伦, 因为把我们
作调和替换的那些互不相交的闭球连续地收缩便给出一个显式同伦. 这样便如
所述得到 $\gamma\in\Omega_{\tilde{\gamma}}$.

对每个满足 $\mathrm{E}(\tilde{\gamma}(\cdot,t))\geq W/2$ 的 $t$, 球族的
性质 (2) 给出某个 $j(t)$, 使得在
$\frac{r_j(t)}2\mathcal{B}_{j(t)}$ 上对 $\tilde{\gamma}(\cdot,t)$ 作
调和替换使能量至少下降
$\frac{e_{\tilde{\gamma},\epsilon_1/8}(t)}{8}$. 于是, 即使在最坏的情形,
即 $r_j(t)\mathcal{B}_{j(t)}$ 是我们在 $t$ 处作替换的第二个球族,
引理~\ref{lem:5.11} 中的 \eqref{eq:5.37} 也给出
\begin{equation}
\mathrm{E}\bigl(\tilde{\gamma}(\cdot,t)\bigr)
-\mathrm{E}\bigl(\gamma(\cdot,t)\bigr)
\geq\kappa\Bigl(\frac{e_{\tilde{\gamma},\epsilon_1/8}(t)}{8}\Bigr)^2.
\label{eq:5.112}
\end{equation}

为建立 $(B_\Psi)$, 设 $\mathcal{B}$ 是 $\bS^2$ 中一族有限个互不相交的
闭球, 使 $\gamma(\cdot,t)$ 在其上的能量至多 $\epsilon_1/12$. 我们可以
假设每个 $\gamma^k(\cdot,t)$ 在 $\mathcal{B}$ 上的能量至多 $\epsilon_1/8$,
否则定理~\ref{thm:5.8} 便给出
$\mathrm{E}(\tilde{\gamma}(\cdot,t))-\mathrm{E}(\gamma(\cdot,t))$ 的正
下界. 于是可以两次应用引理~\ref{lem:5.11} 中的 \eqref{eq:5.38}(先用
$\mu=1/8$, 再用 $\mu=1/4$), 得到
\begin{equation}
\begin{aligned}
\mathrm{E}\bigl(\gamma(\cdot,t)\bigr)
-\mathrm{E}\Bigl[H\bigl(\gamma(\cdot,t),\tfrac18\mathcal{B}\bigr)\Bigr]
&\leq\mathrm{E}\bigl(\tilde{\gamma}(\cdot,t)\bigr)
-\mathrm{E}\Bigl[H\bigl(\tilde{\gamma}(\cdot,t),\tfrac12\mathcal{B}\bigr)\Bigr]\\
&\quad+\frac{2}{\kappa}\Bigl(\mathrm{E}\bigl(\tilde{\gamma}(\cdot,t)\bigr)
-\mathrm{E}\bigl(\gamma(\cdot,t)\bigr)\Bigr)^{1/2}\\
&\leq e_{\tilde{\gamma},\frac{\epsilon_1}{8}}(t)
+\frac{2}{\kappa}\Bigl(\mathrm{E}\bigl(\tilde{\gamma}(\cdot,t)\bigr)
-\mathrm{E}\bigl(\gamma(\cdot,t)\bigr)\Bigr)^{1/2}.
\end{aligned}
\label{eq:5.113}
\end{equation}
把 \eqref{eq:5.112} 与 \eqref{eq:5.113} 同定理~\ref{thm:5.8} 结合起来
即得 $(B_\Psi)$, 从而完成证明.
\end{proof}

\mtpara{好扫掠的存在性}
下一个结果给出一列好扫掠的存在性.

\begin{theorem}
\label{thm:5.26}
(Colding--Minicozzi, \cite{CM27}) 给定 $M$ 上的度量 $g$ 以及代表
$\pi_3(M)$ 中非平凡类的映射 $\beta\in\Omega$, 存在一列扫掠
$\gamma^j\in\Omega_\beta$, 满足
$\max_{s\in[0,1]}\mathrm{E}(\gamma_s^j)\to W(g)$, 并且给定
$\epsilon>0$, 存在 $\bar{j}$ 与 $\delta>0$, 使得若 $j>\bar{j}$ 且
\begin{equation}
\Area\bigl(\gamma^j(\cdot,s)\bigr)>W(g)-\delta,
\label{eq:5.114}
\end{equation}
则存在有限多个调和映射 $u_i:\bS^2\to M$, 使
\begin{equation}
d_V\Bigl(\gamma^j(\cdot,s),\bigcup_i\{u_i\}\Bigr)<\epsilon.
\label{eq:5.115}
\end{equation}
\end{theorem}

在 \eqref{eq:5.115} 中, 我们把每个映射 $u_i$ 与对应于二元组
$(u_i,\bS^2)$ 的变分流形等同起来, 再取这些 $\bS^2$ 的不交并, 得到
$\bigcup_i\{u_i\}$. \eqref{eq:5.115} 中的距离 $d_V$ 是第 3 章
\eqref{eq:3.36} 中定义的变分流形距离.

定理~\ref{thm:5.26} 将在 \S 5.6 中证明.

定理~\ref{thm:5.26} 的一个直接推论是: 若 $s_j$ 是任一序列, 使得
$\Area(\gamma^j(\cdot,s_j))$ 当 $j\to\infty$ 时收敛到宽度 $W(g)$, 则
$\gamma^j(\cdot,s_j)$ 的某个子列收敛到一族从 $\bS^2$ 到 $M$ 的调和映射.
特别地, 这些映射的面积之和恰为 $W(g)$, 而由于它们自动是共形的, 其能量
之和也为 $W(g)$. 从 $\bS^2$ 到 $M$ 至少存在一个非平凡调和映射这一点最早
在 \cite{SaUh} 中证明, 但他们允许极限中有能量损失; 另见 \cite{St}. 在各种
情形下排除这种可能的能量损失被称为``能量等式'', 它可以相当微妙. 这种能量
损失由 Siu 与 Yau 排除, 其中还用了 Meeks 与 Yau 的论证(见 \cite{ScYa2}
第 VIII 章). 后来 Jost 在 \cite{Jo} 的定理 4.2.1 中也证明了这一点, 它给出
至少一个极小极大序列收敛到一族调和映射. \cite{Jo} 中的收敛是在另一种拓扑
下, 正如我们将看到的, 它蕴含变分流形收敛. 关于能量等式的进一步结果见
\cite{Pk}、\cite{DiTi}、\cite{Lt}、\cite{LfWc1} 与 \cite{LfWc2}.

\mtpara{由 $\Omega$ 上的能量下降映射构造好扫掠}
有了定理~\ref{thm:5.23} 与命题~\ref{prop:5.12}, 我们来证明
定理~\ref{thm:5.26}. 记 $\mathcal{G}^{W+1}$ 为从 $\bS^2$ 到 $M$ 的、
能量之和至多 $W+1$ 的调和映射族之集.

\begin{proof}[定理~\ref{thm:5.26}的证明]
选取一列映射 $\tilde{\gamma}^j\in\Omega_\beta$, 使得
\begin{equation}
\max_{t\in[0,1]}\mathrm{E}\bigl(\tilde{\gamma}^j(\cdot,t)\bigr)
<W+\frac1j,
\label{eq:5.116}
\end{equation}
并且 $\tilde{\gamma}^j(\cdot,t)$ 除为常值映射外不是调和的.\footnote{为此,
首先用引理~\ref{lem:5.21}(光滑映射的稠密性), 取
$\tilde{\gamma}_1^j\in\Omega_\beta$ 使
$t\mapsto\tilde{\gamma}_1^j(\cdot,t)$ 从 $[0,1]$ 到 $C^2$ 连续, 且
$\max_{t\in[0,1]}\mathrm{E}(\tilde{\gamma}_1^j(\cdot,t))
<W+\frac1{2j}$. 用球极投影可把 $\tilde{\gamma}_1^j(\cdot,t)$ 视为
$\R^2$ 上的映射. 现在固定 $j$. 连续性在 $C^2$ 中给出一个界
$\sup_{t\in[0,1]}\sup_{B_1}|\nabla\tilde{\gamma}_1^j(\cdot,t)|^2\leq C$,
其中 $C$ 可选取. 取 $R>0$ 使 $4\pi CR^2\leq 1/(2j)$. 定义映射
$\Phi:\R^2\to\R^2$, 在极坐标下为: $\Phi(r,\theta)=(2r,\theta)$ 当
$r<R/2$; $\Phi(r,\theta)=(R,\theta)$ 当 $R/2\leq r\leq R$; 以及
$\Phi(r,\theta)=(r,\theta)$ 当 $r>R$. 注意 $\Phi$ 与恒等映射同伦,
其同伦在圆盘 $B_R\setminus B_{R/2}$ 之外为恒等, 而在
$B_R\setminus B_{R/2}$ 上有 $|\partial_r\Phi|=0$ 与 $|d\Phi|\leq 2$.
于是 $\gamma_2^j(\cdot,t)=\tilde{\gamma}_1^j(\cdot,t)\circ\Phi$ 属于
$\Omega_\beta$, 满足 \eqref{eq:5.116}, 且在 $B_R\setminus B_{R/2}$ 上
有 $\partial_r\gamma_2^j(\cdot,t)=0$. 由于 $\bS^2$ 上的调和映射是共形的
(\lemref{lem:4.25}), 任何调和映射 $\tilde{\gamma}^j(\cdot,t)$ 在
$B_R\setminus B_{R/2}$ 上必为常值, 从而由唯一延拓(定理~1.1 于 \cite{Sj})
在 $\bS^2$ 上为常值.} 可以假设 $W>0$, 否则
$\Area(\tilde{\gamma}^j(\cdot,t))\leq\mathrm{E}(\tilde{\gamma}^j(\cdot,t))
\to 0$, 定理便显然成立.

对 $\tilde{\gamma}^j$ 应用定理~\ref{thm:5.23}, 得到一列
$\gamma^j\in\Omega_\beta$, 其中每个 $\gamma^j(\cdot,t)$ 的能量不超过
$\tilde{\gamma}^j(\cdot,t)$ 的能量. 我们用反证法证明 $\gamma^j$ 具有
所需性质. 因此设存在 $j_k\to\infty$ 与 $s_k\in[0,1]$, 使得
\begin{equation*}
d_V\bigl(\gamma^{j_k}(\cdot,s_k),\mathcal{G}^{W+1}\bigr)\geq\epsilon>0,
\end{equation*}
\begin{equation*}
\Area\bigl(\gamma^{j_k}(\cdot,s_k)\bigr)>W-1/k.
\end{equation*}
于是由 \eqref{eq:5.116} 以及 $\mathrm{E}(\cdot)\geq\Area(\cdot)$ 得到
\begin{equation}
\begin{aligned}
\mathrm{E}\bigl(\tilde{\gamma}^{j_k}(\cdot,s_k)\bigr)
-\mathrm{E}\bigl(\gamma^{j_k}(\cdot,s_k)\bigr)
&\leq\mathrm{E}\bigl(\tilde{\gamma}^{j_k}(\cdot,s_k)\bigr)
-\Area\bigl(\gamma^{j_k}(\cdot,s_k)\bigr)\\
&\leq 1/k+1/j_k\to 0,
\end{aligned}
\label{eq:5.117}
\end{equation}
类似地,
\begin{equation*}
\mathrm{E}\bigl(\gamma^{j_k}(\cdot,s_k)\bigr)
-\Area\bigl(\gamma^{j_k}(\cdot,s_k)\bigr)\to 0.
\end{equation*}
把 \eqref{eq:5.117} 用于定理~\ref{thm:5.23}, 得到

\medskip
\noindent (B) 若 $\mathcal{B}$ 是 $\bS^2$ 中任一由互不相交的闭球构成
的族, 且
$\int_{\bigcup_{\mathcal{B}}B}|\nabla\gamma^{j_k}(\cdot,s_k)|^2
<\epsilon_0$, 而 $v:\bigcup_{\mathcal{B}}\frac18B\to M$ 是在
$\bigcup_{\mathcal{B}}\frac18\partial B$ 上等于
$\gamma^{j_k}(\cdot,s_k)$ 的极小化能量映射, 则
\begin{equation}
\int_{\bigcup_{\mathcal{B}}\frac18B}
\bigl|\nabla\gamma^{j_k}(\cdot,s_k)-\nabla v\bigr|^2
\leq\Psi(1/k+1/j_k)\to 0.
\label{eq:5.118}
\end{equation}
因此可以应用命题~\ref{prop:5.12}, 得到 $\gamma^{j_k}(\cdot,s_k)$ 的某个
子列变分流形收敛到一族调和映射. 但这与它到
$\mathcal{G}^{W+1}$ 的变分流形距离有正下界相矛盾, 从而完成证明.
\end{proof}

\mtpara{参数空间}
我们本可以用任意紧的有限维拓扑空间 $\mathcal{P}$ 作为扫掠中映射的参数
空间, 只要求映射在 $\partial\mathcal{P}$ 上为常值(或者
$\partial\mathcal{P}=\varnothing$), 而不必用单位区间 $[0,1]$ 并要求映射
起止于点映射. 此时记 $\Omega^{\mathcal{P}}$ 为连续映射
$\sigma:\bS^2\times\mathcal{P}\to M$ 之集, 使得对每个
$t\in\mathcal{P}$, 映射 $\sigma(\cdot,t)$ 属于
$C^0\cap W^{1,2}(\bS^2,M)$, 映射 $t\mapsto\sigma(\cdot,t)$ 是从
$\mathcal{P}$ 到 $C^0\cap W^{1,2}(\bS^2,M)$ 的连续映射, 并且最后
$\sigma$ 把 $\partial\mathcal{P}$ 映为点映射. 给定映射
$\hat{\sigma}\in\Omega^{\mathcal{P}}$, 同伦类
$\Omega_{\hat{\sigma}}^{\mathcal{P}}\subset\Omega^{\mathcal{P}}$ 定义为
$\Omega^{\mathcal{P}}$ 中所有与 $\hat{\sigma}$ 在
$\Omega^{\mathcal{P}}$ 内同伦的映射之集. 最后, 宽度
$W=W(\hat{\sigma})$ 为
\begin{equation*}
\inf_{\sigma\in\Omega_{\hat{\sigma}}^{\mathcal{P}}}
\max_{t\in\mathcal{P}}\mathrm{E}\bigl(\sigma(\cdot,t)\bigr).
\end{equation*}
只需作平凡的改动, 同样的证明便给出这些一般参数空间上的
定理~\ref{thm:5.26}.\footnote{主要改动在引理~\ref{lem:5.25}, 其中
性质 (1) 里重数的界 $2$ 变为 $\dim(\mathcal{P})+1$. 这由覆盖维数的定义
推出; 见 \cite{Mu} 第 302--303 页.}
